%% Generated by Sphinx.
\def\sphinxdocclass{report}
\documentclass[letterpaper,11pt,english]{sphinxmanual}
\ifdefined\pdfpxdimen
   \let\sphinxpxdimen\pdfpxdimen\else\newdimen\sphinxpxdimen
\fi \sphinxpxdimen=.75bp\relax
\ifdefined\pdfimageresolution
    \pdfimageresolution= \numexpr \dimexpr1in\relax/\sphinxpxdimen\relax
\fi
\newdimen\sphinxremdimen\sphinxremdimen = 11pt
%% let collapsible pdf bookmarks panel have high depth per default
\PassOptionsToPackage{bookmarksdepth=5}{hyperref}
%% turn off hyperref patch of \index as sphinx.xdy xindy module takes care of
%% suitable \hyperpage mark-up, working around hyperref-xindy incompatibility
\PassOptionsToPackage{hyperindex=false}{hyperref}
%% memoir class requires extra handling
\makeatletter\@ifclassloaded{memoir}
{\ifdefined\memhyperindexfalse\memhyperindexfalse\fi}{}\makeatother

\PassOptionsToPackage{booktabs}{sphinx}
\PassOptionsToPackage{colorrows}{sphinx}

\PassOptionsToPackage{warn}{textcomp}

\catcode`^^^^00a0\active\protected\def^^^^00a0{\leavevmode\nobreak\ }
\usepackage{cmap}
\usepackage{fontspec}
\defaultfontfeatures[\rmfamily,\sffamily,\ttfamily]{}
\usepackage{amsmath,amssymb,amstext}
\usepackage{polyglossia}
\setmainlanguage{english}



\setmainfont{FreeSerif}[
  Extension      = .otf,
  UprightFont    = *,
  ItalicFont     = *Italic,
  BoldFont       = *Bold,
  BoldItalicFont = *BoldItalic
]
\setsansfont{FreeSans}[
  Extension      = .otf,
  UprightFont    = *,
  ItalicFont     = *Oblique,
  BoldFont       = *Bold,
  BoldItalicFont = *BoldOblique,
]
\setmonofont{FreeMono}[Scale=0.9,
  Extension      = .otf,
  UprightFont    = *,
  ItalicFont     = *Oblique,
  BoldFont       = *Bold,
  BoldItalicFont = *BoldOblique,
]



\usepackage[Bjarne]{fncychap}
\usepackage{sphinx}

\fvset{fontsize=auto}
\usepackage{geometry}


% Include hyperref last.
\usepackage{hyperref}
% Fix anchor placement for figures with captions.
\usepackage{hypcap}% it must be loaded after hyperref.
% Set up styles of URL: it should be placed after hyperref.
\urlstyle{same}

\addto\captionsenglish{\renewcommand{\contentsname}{sd-boot Documentation}}

\usepackage{sphinxmessages}
\setcounter{tocdepth}{1}


    \usepackage{parskip}
    \usepackage{fontspec}

    % Fix the 11pt headheight layout warnings
    \setlength{\headheight}{14pt}
    \addtolength{\topmargin}{-2pt}

    % Strip vertical spaces between items
    \usepackage{enumitem}
    \setlist[itemize]{
        noitemsep, 
        topsep=6pt,             % \parskip,       % 0pt, 
        parsep=0pt, 
        partopsep=0pt,
        after=\vspace{0pt}
        }
    \setlist[enumerate]{
        noitemsep, 
        topsep=6pt,             % 0pt, 
        parsep=0pt, 
        partopsep=0pt,
        after=\vspace{0pt}
        }

    \usepackage{newunicodechar}
    \newunicodechar{␣}{\textvisiblespace}
    \tracinglostchars=0

    % Body sans serif 
    % \setmainfont{TeX Gyre Heros}[
    % \setmainfont{IBM Plex Sans}[
    \setmainfont{Source Sans 3}[
        Ligatures=TeX,
        Scale=0.92
    ]

    % Section Headers (Modern Helvetica equivalent)
    % \setsansfont{TeX Gyre Heros}[
    % \setsansfont{IBM Plex Sans}[
    \setsansfont{Source Sans 3}[
        Ligatures=TeX,
        Scale=0.92
    ]

    % Verbatim/Inline (Monospace Fira)
    \setmonofont{Fira Mono}[
        Scale=0.88
    ]
    

\title{sd-boot Documentation }
\date{Sep 28, 2026}
\release{6.5.1
}
\author{Gene C}
\newcommand{\sphinxlogo}{\vbox{}}
\renewcommand{\releasename}{Release}
\makeindex
\begin{document}

\pagestyle{empty}
\sphinxmaketitle
\pagestyle{plain}
\sphinxtableofcontents
\pagestyle{normal}
\phantomsection\label{\detokenize{index::doc}}


\sphinxstepscope


\chapter{sd\sphinxhyphen{}boot}
\label{\detokenize{Readme:sd-boot}}\label{\detokenize{Readme:read-me}}\label{\detokenize{Readme::doc}}

\section{Overview}
\label{\detokenize{Readme:overview}}
\sphinxAtStartPar
\sphinxstylestrong{sd\sphinxhyphen{}boot} installs linux kernels and efi programs to the \sphinxstyleemphasis{boot} partition using systemd’s \sphinxstyleemphasis{kernel\sphinxhyphen{}install}.
This is a robust and systematic way to install or remove bootable software to or
from the \sphinxstyleemphasis{boot} partition.

\sphinxAtStartPar
\sphinxstyleemphasis{kernel\sphinxhyphen{}install} uses \sphinxstyleemphasis{\$BOOT} to refer to this partition or it’s mount point.
It is usually mounted on either \sphinxstyleemphasis{/efi} or \sphinxstyleemphasis{/boot}.
In the past it \sphinxstyleemphasis{/boot/efi} has also been used, but mounting the ESP under /boot is now strongly discouraged.

\sphinxAtStartPar
The boot process requires access to the ESP partition
and there are two distinct recommendations.

\sphinxAtStartPar
Mount the ESP partition:
\begin{itemize}
\item {} 
\sphinxAtStartPar
onto \sphinxstyleemphasis{/efi} and an XBOOTLDR partition, if any, onto \sphinxstyleemphasis{/boot}

\item {} 
\sphinxAtStartPar
onto \sphinxstyleemphasis{/boot} or \sphinxstyleemphasis{/efi} if there is an XBOOTLDR partition and mount that on \sphinxstyleemphasis{/boot}.

\end{itemize}

\sphinxAtStartPar
The first is recommended by kernel\sphinxhyphen{}install and the second by the
\sphinxhref{https://uapi-group.org/specifications/specs/boot\_loader\_specification}{UAPI Group Specifications: Mount Points}.

\sphinxAtStartPar
sd\sphinxhyphen{}boot and kernel\sphinxhyphen{}install both happily work with either approach.

\sphinxAtStartPar
Our preference follows kernel\sphinxhyphen{}install; mount the ESP on \sphinxstyleemphasis{/efi} and an
XBOOTLDR partition on \sphinxstyleemphasis{/boot}. This is clean, simple and unambiguous.

\sphinxAtStartPar
We designate the ESP mount point location as \sphinxstyleemphasis{<EFI>}, which is either \sphinxstyleemphasis{/efi} or \sphinxstyleemphasis{/boot} as
appropriate. We also refer to the directory where kernels are put as \sphinxstyleemphasis{<BOOT>}.
If you mount the ESP onto \sphinxstyleemphasis{/boot} then both \sphinxstyleemphasis{<EFI>} and \sphinxstyleemphasis{<BOOT>} refer to \sphinxstyleemphasis{/boot}.
If the ESP is mounted on \sphinxstyleemphasis{/efi} and you have an XBOOTLDR partition mounted on \sphinxstyleemphasis{/boot}
then they refer to those (obviously).

\sphinxAtStartPar
While we recommend using the first approach for these mount points, we must also support
any systems where the ESP is mounted on \sphinxstyleemphasis{/boot}. This guarantees that all systems
continue to work without any change.

\sphinxAtStartPar
sd\sphinxhyphen{}boot provides:
\begin{itemize}
\item {} 
\sphinxAtStartPar
pacman alpm hooks

\sphinxAtStartPar
These ensure that kernels and efi\sphinxhyphen{}tools  are installed or removed from the boot partition when
a kernel package is installed, updated or removed.

\item {} 
\sphinxAtStartPar
installers that are triggered from ALPM hooks.

\sphinxAtStartPar
These do the actual work.

\item {} 
\sphinxAtStartPar
command line tools

\sphinxAtStartPar
Give administrators a simple way to directly (re\sphinxhyphen{})install or remove kernels from the boot partition.

\item {} 
\sphinxAtStartPar
kernel\sphinxhyphen{}install plugin(s)

\sphinxAtStartPar
Used for loader entries for non\sphinxhyphen{}UKI kernels and efi tools.

\end{itemize}

\sphinxAtStartPar
\sphinxstyleemphasis{Installing} a kernel (or efi tool) means placing the appropriate image in \sphinxstyleemphasis{\$BOOT} partition so
that the image can be booted.

\sphinxAtStartPar
Any kernel packages must therefore place kernel images in the \sphinxstyleemphasis{/usr} partition only.

\sphinxAtStartPar
The same applies for any bootable efi tools such as \sphinxstyleemphasis{edk2\sphinxhyphen{}shell} and \sphinxstyleemphasis{memtest86+}.


\subsection{Documentation}
\label{\detokenize{Readme:documentation}}
\sphinxAtStartPar
The manual provides detailed information about using sd\sphinxhyphen{}boot.
It is available in both HTML and PDF formats.
Both are installed under \sphinxstyleemphasis{/usr/share/sd\sphinxhyphen{}boot/docs}

\sphinxAtStartPar
Each executable comes with a man page:
\begin{itemize}
\item {} 
\sphinxAtStartPar
sd\sphinxhyphen{}boot\sphinxhyphen{}efifs\sphinxhyphen{}update

\item {} 
\sphinxAtStartPar
sd\sphinxhyphen{}boot\sphinxhyphen{}efi\sphinxhyphen{}tool\sphinxhyphen{}update

\item {} 
\sphinxAtStartPar
sd\sphinxhyphen{}boot\sphinxhyphen{}find\sphinxhyphen{}boot\sphinxhyphen{}mounts

\item {} 
\sphinxAtStartPar
sd\sphinxhyphen{}boot\sphinxhyphen{}kernel\sphinxhyphen{}update

\end{itemize}

\sphinxAtStartPar
For convenience, the manual includes all the man pages.


\subsection{Coding Standards}
\label{\detokenize{Readme:coding-standards}}
\sphinxAtStartPar
All code should be carefully reviewed and checked. This is particularly true for
system utilities where the importance of avoiding potential problems is high.

\sphinxAtStartPar
Placing trust in software requires care to ensure that it not only performs
the required tasks but that the code is robust, handles exceptions gracefully and does
not suffer from potential memory related issues. And of course no bugs!
Sure, having zero bugs ever is tricky. While we cannot guarantee perfection,
there are things we can do to minimize them and bolster confidence in the code.

\sphinxAtStartPar
To help with this, sd\sphinxhyphen{}boot:
\begin{itemize}
\item {} 
\sphinxAtStartPar
has undergone functional testing

\item {} 
\sphinxAtStartPar
has undergone static analysis with tools such as clang\sphinxhyphen{}tidy

\item {} 
\sphinxAtStartPar
has undergone compiler analysis using tools such as \sphinxhyphen{}fanalyzer

\item {} 
\sphinxAtStartPar
is completely clean when run under valgrind

\item {} 
\sphinxAtStartPar
has undergone human code review

\item {} 
\sphinxAtStartPar
has underfone additional AI reviews.
\begin{itemize}
\item {} 
\sphinxAtStartPar
Assisted\sphinxhyphen{}by: Claude (Anthropic) <\sphinxurl{https://claude.ai}>

\end{itemize}

\end{itemize}


\section{Quick Start}
\label{\detokenize{Readme:quick-start}}
\sphinxAtStartPar
There are two sets of configuration files
\begin{itemize}
\item {} 
\sphinxAtStartPar
kernel\sphinxhyphen{}install configuration

\item {} 
\sphinxAtStartPar
sd\sphinxhyphen{}boot configuration

\end{itemize}

\sphinxAtStartPar
Both follow the Linux standard priority order and \sphinxstyleemphasis{drop\sphinxhyphen{}in} files that over\sphinxhyphen{}ride
settings in lower priority files. In rough terms:
\begin{itemize}
\item {} 
\sphinxAtStartPar
Files within a directory are prioritized alpha\sphinxhyphen{}numerically.

\item {} 
\sphinxAtStartPar
Files in \sphinxstyleemphasis{/etc} over\sphinxhyphen{}ride those in \sphinxstyleemphasis{/usr}

\item {} 
\sphinxAtStartPar
System provided files reside in \sphinxstyleemphasis{/usr}

\item {} 
\sphinxAtStartPar
Administration settings reside in \sphinxstyleemphasis{/etc}

\item {} 
\sphinxAtStartPar
Drop\sphinxhyphen{}in files in directory \sphinxstyleemphasis{xxx.d} may override partial settings.

\end{itemize}

\sphinxAtStartPar
See \sphinxstyleemphasis{man 5 systemd.unit} for additional information about drop\sphinxhyphen{}in configuration files.


\subsection{kernel\sphinxhyphen{}install}
\label{\detokenize{Readme:kernel-install}}
\sphinxAtStartPar
Please see the kernel\sphinxhyphen{}install documentation for full details.
sd\sphinxhyphen{}boot provides:
\begin{itemize}
\item {} 
\sphinxAtStartPar
/usr/lib/kernel/install.conf.d/010\sphinxhyphen{}sd\sphinxhyphen{}boot\sphinxhyphen{}install.conf

\end{itemize}

\sphinxAtStartPar
which sets the following:

\begin{sphinxVerbatim}[commandchars=\\\{\}]
\PYG{n}{layout}\PYG{o}{=}\PYG{n}{uki}
\PYG{n}{initrd\PYGZus{}generator}\PYG{o}{=}\PYG{n}{dracut}
\PYG{n}{uki\PYGZus{}generator}\PYG{o}{=}\PYG{n}{ukify}
\end{sphinxVerbatim}

\sphinxAtStartPar
These may be overridden in:
\begin{itemize}
\item {} 
\sphinxAtStartPar
/etc/kernel/install.conf

\item {} 
\sphinxAtStartPar
/etc/kernel/install.conf.d/xxx.conf

\end{itemize}

\sphinxAtStartPar
We recommend sticking with UKI and dracut.


\subsection{sd\sphinxhyphen{}boot}
\label{\detokenize{Readme:id1}}
\sphinxAtStartPar
All the configuration files are located in \sphinxstyleemphasis{/etc/sd\sphinxhyphen{}boot}:
\begin{itemize}
\item {} 
\sphinxAtStartPar
\sphinxstylestrong{/etc/sd\sphinxhyphen{}boot/config.yaml}

\sphinxAtStartPar
Settings:

\begin{sphinxVerbatim}[commandchars=\\\{\}]
\PYG{n}{verb}\PYG{p}{:} \PYG{l+m+mi}{1}
\PYG{n}{skip\PYGZus{}kernel\PYGZus{}plugins}\PYG{p}{:} \PYG{p}{[}\PYG{p}{]}
\end{sphinxVerbatim}

\item {} 
\sphinxAtStartPar
\sphinxstylestrong{/etc/sd\sphinxhyphen{}boot/efi\sphinxhyphen{}tool.packages}

\sphinxAtStartPar
A list of all package names of efi tools to be managed by sd\sphinxhyphen{}boot:

\begin{sphinxVerbatim}[commandchars=\\\{\}]
\PYG{n}{memtest86\PYGZus{}64}\PYG{o}{\PYGZhy{}}\PYG{n}{git}
\PYG{n}{edk2}\PYG{o}{\PYGZhy{}}\PYG{n}{shell}
\end{sphinxVerbatim}

\item {} 
\sphinxAtStartPar
\sphinxstylestrong{/etc/sd\sphinxhyphen{}boot/kernel.packages}

\sphinxAtStartPar
A list of kernel packages managed by sd\sphinxhyphen{}boot:

\begin{sphinxVerbatim}[commandchars=\\\{\}]
\PYG{n}{linux}\PYG{o}{\PYGZhy{}}\PYG{n}{custom}
\PYG{n}{linux}\PYG{o}{\PYGZhy{}}\PYG{n+nb}{next}
\PYG{n}{linux}\PYG{o}{\PYGZhy{}}\PYG{n}{test}
\PYG{n}{linux}\PYG{o}{\PYGZhy{}}\PYG{n}{stable}
\end{sphinxVerbatim}

\end{itemize}

\sphinxAtStartPar
Unlike kernels whihc are all installed in a fixed location,
efi tools must provide an \sphinxstyleemphasis{<package\sphinxhyphen{}name>.image} file which has the
full path to the \sphinxstyleemphasis{.efi} program.

\sphinxAtStartPar
These are included in sd\sphinxhyphen{}boot package:
\begin{itemize}
\item {} 
\sphinxAtStartPar
/etc/sd\sphinxhyphen{}boot/edk2\sphinxhyphen{}shell.image

\item {} 
\sphinxAtStartPar
/etc/sd\sphinxhyphen{}boot/memtest86\_64\sphinxhyphen{}git.image

\end{itemize}

\sphinxAtStartPar
Comment lines starting with \sphinxstylestrong{\#} are permitted in all config files.


\subsection{dracut}
\label{\detokenize{Readme:dracut}}
\sphinxAtStartPar
sd\sphinxhyphen{}boot provides a default dracut configuration in /usr/lib/dracut/dracut.conf.d/010\sphinxhyphen{}sd\sphinxhyphen{}boot\sphinxhyphen{}dracut.conf
which may be overridden in:
\begin{itemize}
\item {} 
\sphinxAtStartPar
\sphinxstylestrong{/etc/dracut.conf.d/010\sphinxhyphen{}sd\sphinxhyphen{}boot\sphinxhyphen{}dracut.conf}

\end{itemize}

\sphinxstepscope


\chapter{Configuration}
\label{\detokenize{Configuration:configuration}}\label{\detokenize{Configuration::doc}}
\sphinxAtStartPar
In addition to the files kernel\sphinxhyphen{}install itself uses, such as
those in the \sphinxstyleemphasis{/etc/kernel/} directory, the \sphinxstyleemphasis{sd\sphinxhyphen{}boot} configuration files
reside in \sphinxstyleemphasis{/etc/sd\sphinxhyphen{}boot/}.

\sphinxAtStartPar
The file \sphinxstyleemphasis{/etc/sd\sphinxhyphen{}boot/kernel.packages} lists the kernel package names
sd\sphinxhyphen{}boot is permitted to install.

\sphinxAtStartPar
By default, the Arch kernel is not included in the list, but of course sd\sphinxhyphen{}boot can manage that too.
The section {\hyperref[\detokenize{Arch-Kernels:arch-kernels}]{\sphinxcrossref{\DUrole{std}{\DUrole{std-ref}{Installing Arch Kernels With sd\sphinxhyphen{}boot}}}}} explains how to use sd\sphinxhyphen{}boot for the Arch kernel.

\sphinxAtStartPar
Similarly, \sphinxstyleemphasis{/etc/sd\sphinxhyphen{}boot/efi\sphinxhyphen{}tools.packages} lists the efi tool packages that
sd\sphinxhyphen{}boot is permitted to install.

\sphinxAtStartPar
Package names here are the standard Arch package names.

\sphinxAtStartPar
Please ensure any package to be managed by sd\sphinxhyphen{}boot is listed appropriately.

\sphinxAtStartPar
By default kernel\sphinxhyphen{}install sets kernel boot options using in order:

\begin{sphinxVerbatim}[commandchars=\\\{\}]
\PYG{o}{/}\PYG{n}{etc}\PYG{o}{/}\PYG{n}{kernel}\PYG{o}{/}\PYG{n}{cmdline}
\PYG{o}{/}\PYG{n}{usr}\PYG{o}{/}\PYG{n}{lib}\PYG{o}{/}\PYG{n}{kernel}\PYG{o}{/}\PYG{n}{cmdline}
\PYG{o}{/}\PYG{n}{proc}\PYG{o}{/}\PYG{n}{cmdline}
\end{sphinxVerbatim}

\sphinxAtStartPar
Putting kernel command line options into \sphinxstyleemphasis{/etc/kernel/cmdline} will then over\sphinxhyphen{}ride the
default. If no file provides the kernel options, then The default is to use /proc/cmdline
which proivdes the kernel command line option of the currently booted kernel.

\sphinxAtStartPar
On initial install sd\sphinxhyphen{}boot provides
\sphinxstyleemphasis{/etc/kernel/install.conf} which sets the layout to \sphinxstyleemphasis{uki}, the initrd generator to \sphinxstyleemphasis{dracut}
and the uki generator to \sphinxstyleemphasis{ukify}.

\sphinxAtStartPar
The \sphinxhref{https://gitlab.archlinux.org/archlinux/rfcs/-/merge\_requests/66\#note\_452083}{RFC 66 Discussion}
about changing to use \sphinxstyleemphasis{kernel\sphinxhyphen{}install} triggered me to migrate my own kernels and
see how it works in practice. I found it works really well and am sharing this
with the community in case its helpful to others.

\sphinxAtStartPar
The original version of \sphinxstyleemphasis{sd\sphinxhyphen{}boot} was written in bash, but the current version
coded in \sphinxstyleemphasis{C} has since replaced it.


\section{Getting Started}
\label{\detokenize{Configuration:getting-started}}
\sphinxAtStartPar
For a kernel to be managed by \sphinxstyleemphasis{sd\sphinxhyphen{}boot} the kernel package should install to:

\begin{sphinxVerbatim}[commandchars=\\\{\}]
\PYG{o}{/}\PYG{n}{usr}\PYG{o}{/}\PYG{n}{lib}\PYG{o}{/}\PYG{n}{modules}\PYG{o}{/}\PYG{o}{\PYGZlt{}}\PYG{n}{kernel}\PYG{o}{\PYGZhy{}}\PYG{n}{version}\PYG{o}{\PYGZgt{}}
\end{sphinxVerbatim}

\sphinxAtStartPar
with the kernel image at:

\begin{sphinxVerbatim}[commandchars=\\\{\}]
\PYG{o}{/}\PYG{n}{usr}\PYG{o}{/}\PYG{n}{lib}\PYG{o}{/}\PYG{n}{modules}\PYG{o}{/}\PYG{o}{\PYGZlt{}}\PYG{n}{kernel}\PYG{o}{\PYGZhy{}}\PYG{n}{version}\PYG{o}{\PYGZgt{}}\PYG{o}{/}\PYG{n}{vmlinuz}
\end{sphinxVerbatim}

\sphinxAtStartPar
A kernel package must \sphinxstylestrong{NOT} install the kernel or any initrd into \sphinxstyleemphasis{<EFI>} or \sphinxstyleemphasis{<BOOT>}.
It should not even create an initrd. Installing anything to these \sphinxstyleemphasis{\$BOOT} directories
must be left to sd\sphinxhyphen{}boot. This applies equally to any efi tools.

\sphinxAtStartPar
Once the package is installed, then grant sd\sphinxhyphen{}boot permission to manage
those kernels and efi\sphinxhyphen{}tools by listing them in the files:

\begin{sphinxVerbatim}[commandchars=\\\{\}]
\PYG{o}{/}\PYG{n}{etc}\PYG{o}{/}\PYG{n}{sd}\PYG{o}{\PYGZhy{}}\PYG{n}{boot}\PYG{o}{/}\PYG{n}{kernel}\PYG{o}{.}\PYG{n}{packages}
\PYG{o}{/}\PYG{n}{etc}\PYG{o}{/}\PYG{n}{sd}\PYG{o}{\PYGZhy{}}\PYG{n}{boot}\PYG{o}{/}\PYG{n}{efi}\PYG{o}{\PYGZhy{}}\PYG{n}{tools}\PYG{o}{.}\PYG{n}{packages}
\end{sphinxVerbatim}

\sphinxAtStartPar
On the next update (or re\sphinxhyphen{}install or remove), if the kernel or efi\sphinxhyphen{}tool is designated
to be managed by \sphinxstyleemphasis{sd\sphinxhyphen{}boot} then it will handle installing it into \$BOOT.

\sphinxAtStartPar
You can (re\sphinxhyphen{})install a kernel at any time at the command line using:

\begin{sphinxVerbatim}[commandchars=\\\{\}]
\PYG{n}{sd}\PYG{o}{\PYGZhy{}}\PYG{n}{boot}\PYG{o}{\PYGZhy{}}\PYG{n}{kernel}\PYG{o}{\PYGZhy{}}\PYG{n}{install} \PYG{o}{\PYGZlt{}}\PYG{n}{kernel} \PYG{n}{package} \PYG{n}{name}\PYG{o}{\PYGZgt{}}
\end{sphinxVerbatim}

\sphinxAtStartPar
or trigger a pacman refresh with :

\begin{sphinxVerbatim}[commandchars=\\\{\}]
touch\PYG{+w}{ }/usr/lib/modules/\PYGZlt{}kernel\PYGZhy{}version\PYGZgt{}/vmlinuz
pacman\PYG{+w}{ }\PYGZhy{}Syu
\end{sphinxVerbatim}

\sphinxAtStartPar
sd\sphinxhyphen{}boot relies on each kernel package providing a file in the kernel module
directory that contains that kernel package name:

\begin{sphinxVerbatim}[commandchars=\\\{\}]
\PYG{o}{/}\PYG{n}{usr}\PYG{o}{/}\PYG{n}{lib}\PYG{o}{/}\PYG{n}{modules}\PYG{o}{/}\PYG{o}{\PYGZlt{}}\PYG{n}{kernel}\PYG{o}{\PYGZhy{}}\PYG{n}{version}\PYG{o}{\PYGZgt{}}\PYG{o}{/}\PYG{n}{pkgbase}
\end{sphinxVerbatim}

\sphinxAtStartPar
or:

\begin{sphinxVerbatim}[commandchars=\\\{\}]
\PYG{o}{/}\PYG{n}{usr}\PYG{o}{/}\PYG{n}{lib}\PYG{o}{/}\PYG{n}{modules}\PYG{o}{/}\PYG{o}{\PYGZlt{}}\PYG{n}{kernel}\PYG{o}{\PYGZhy{}}\PYG{n}{version}\PYG{o}{\PYGZgt{}}\PYG{o}{/}\PYG{n}{pkgbase}\PYG{o}{\PYGZhy{}}\PYG{n}{sdb}
\end{sphinxVerbatim}

\sphinxAtStartPar
Arch kernels provide \sphinxstyleemphasis{pkgbase} file.
While \sphinxstyleemphasis{sd\sphinxhyphen{}boot} uses either of these files, \sphinxstyleemphasis{pkgbase\sphinxhyphen{}sdb} is preferred
since the Arch kernel install tools act on any kernel using \sphinxstyleemphasis{pkgbase}.
Using a different filename prevents the Arch tools from installing kernels
sd\sphinxhyphen{}boot is already handling.


\section{Command Line Tools}
\label{\detokenize{Configuration:command-line-tools}}
\sphinxAtStartPar
The following tools may be run from the command line:
\begin{itemize}
\item {} 
\sphinxAtStartPar
/usr/bin/sd\sphinxhyphen{}boot\sphinxhyphen{}efi\sphinxhyphen{}tool\sphinxhyphen{}update

\item {} 
\sphinxAtStartPar
/usr/bin/sd\sphinxhyphen{}boot\sphinxhyphen{}kernel\sphinxhyphen{}update

\item {} 
\sphinxAtStartPar
/usr/lib/sd\sphinxhyphen{}boot/sd\sphinxhyphen{}boot\sphinxhyphen{}efifs\sphinxhyphen{}update

\item {} 
\sphinxAtStartPar
/usr/lib/sd\sphinxhyphen{}boot/sd\sphinxhyphen{}boot\sphinxhyphen{}find\sphinxhyphen{}boot\sphinxhyphen{}mounts

\end{itemize}

\sphinxAtStartPar
There is a man page for each of these. For example see \sphinxstyleemphasis{man sd\sphinxhyphen{}boot\sphinxhyphen{}kernel\sphinxhyphen{}update}.

\sphinxAtStartPar
sd\sphinxhyphen{}boot is normally triggered into action from pacman via ALPM hooks. Kernels and efi tools
may also be manually installed or removed from \$BOOT.

\sphinxAtStartPar
For example to install a bootable efi shell, provided by the \sphinxstyleemphasis{adk2\sphinxhyphen{}shell} package:

\begin{sphinxVerbatim}[commandchars=\\\{\}]
sd\PYGZhy{}boot\PYGZhy{}efi\PYGZhy{}tool\PYGZhy{}update\PYG{+w}{ }add\PYG{+w}{ }edk2\PYGZhy{}shell
\end{sphinxVerbatim}

\sphinxAtStartPar
To install a kernel provided by the \sphinxstyleemphasis{linux\sphinxhyphen{}custom} package

\begin{sphinxVerbatim}[commandchars=\\\{\}]
sd\PYGZhy{}boot\PYGZhy{}kernel\PYGZhy{}update\PYG{+w}{ }add\PYG{+w}{ }linux\PYGZhy{}custom
\end{sphinxVerbatim}

\sphinxAtStartPar
Please ensure that any packages to be managed by sd\sphinxhyphen{}boot are listed appropriately in:

\begin{sphinxVerbatim}[commandchars=\\\{\}]
/etc/sd\PYGZhy{}boot/efi\PYGZhy{}tool.packages
/etc/sd\PYGZhy{}boot/kernel.packages
\end{sphinxVerbatim}

\sphinxAtStartPar
To install efi filesystem drivers:

\begin{sphinxVerbatim}[commandchars=\\\{\}]
/usr/lib/sd\PYGZhy{}boot/sd\PYGZhy{}boot\PYGZhy{}efifs\PYGZhy{}update\PYG{+w}{ }add
\end{sphinxVerbatim}

\sphinxAtStartPar
To list ESP and XBOOTLDR partions mount points run (as root):

\begin{sphinxVerbatim}[commandchars=\\\{\}]
/usr/lib/sd\PYGZhy{}boot/sd\PYGZhy{}boot\PYGZhy{}find\PYGZhy{}boot\PYGZhy{}mounts
\end{sphinxVerbatim}

\sphinxAtStartPar
which will list them all and those in use by the current booted system are marked
with an asterisk.

\sphinxAtStartPar
Replace \sphinxstyleemphasis{add} by \sphinxstyleemphasis{remove} to remove them from \$BOOT. Note that \sphinxstyleemphasis{add} and \sphinxstyleemphasis{remove} refers
only to a copy of the image in \$BOOT and does not install or remove the package itself.
That job belongs to pacman.

\sphinxstepscope


\chapter{Installing Arch Kernels With sd\sphinxhyphen{}boot}
\label{\detokenize{Arch-Kernels:installing-arch-kernels-with-sd-boot}}\label{\detokenize{Arch-Kernels:arch-kernels}}\label{\detokenize{Arch-Kernels::doc}}
\sphinxAtStartPar
By default the Archlinux kernel package, \sphinxstyleemphasis{linux}, is not in the list of kernels managed by sd\sphinxhyphen{}boot.

\sphinxAtStartPar
In order for sd\sphinxhyphen{}boot to manage it, there are two steps that must be done.
\begin{itemize}
\item {} 
\sphinxAtStartPar
Add \sphinxstylestrong{linux} to the list of sd\sphinxhyphen{}boot managed kernels:

\begin{sphinxVerbatim}[commandchars=\\\{\}]
\PYG{n}{echo} \PYG{n}{linux} \PYG{o}{\PYGZgt{}\PYGZgt{}} \PYG{o}{/}\PYG{n}{etc}\PYG{o}{/}\PYG{n}{sd}\PYG{o}{\PYGZhy{}}\PYG{n}{boot}\PYG{o}{/}\PYG{n}{kernel}\PYG{o}{.}\PYG{n}{packages}
\end{sphinxVerbatim}

\item {} 
\sphinxAtStartPar
Prevent mkinitcpio from installing a duplicate of the same kernel:

\begin{sphinxVerbatim}[commandchars=\\\{\}]
\PYG{n}{sudo} \PYG{n}{mkdir} \PYG{o}{\PYGZhy{}}\PYG{n}{p} \PYG{o}{/}\PYG{n}{etc}\PYG{o}{/}\PYG{n}{pacman}\PYG{o}{.}\PYG{n}{d}\PYG{o}{/}\PYG{n}{hooks}
\PYG{n}{sudo} \PYG{n}{ln} \PYG{o}{\PYGZhy{}}\PYG{n}{s} \PYG{o}{/}\PYG{n}{dev}\PYG{o}{/}\PYG{n}{null} \PYG{o}{/}\PYG{n}{etc}\PYG{o}{/}\PYG{n}{pacman}\PYG{o}{.}\PYG{n}{d}\PYG{o}{/}\PYG{n}{hooks}\PYG{o}{/}\PYG{l+m+mi}{60}\PYG{o}{\PYGZhy{}}\PYG{n}{mkinitcpio}\PYG{o}{\PYGZhy{}}\PYG{n}{remove}\PYG{o}{.}\PYG{n}{hook}
\PYG{n}{sudo} \PYG{n}{ln} \PYG{o}{\PYGZhy{}}\PYG{n}{s} \PYG{o}{/}\PYG{n}{dev}\PYG{o}{/}\PYG{n}{null} \PYG{o}{/}\PYG{n}{etc}\PYG{o}{/}\PYG{n}{pacman}\PYG{o}{.}\PYG{n}{d}\PYG{o}{/}\PYG{n}{hooks}\PYG{o}{/}\PYG{l+m+mi}{90}\PYG{o}{\PYGZhy{}}\PYG{n}{mkinitcpio}\PYG{o}{\PYGZhy{}}\PYG{n}{install}\PYG{o}{.}\PYG{n}{hook}
\end{sphinxVerbatim}

\end{itemize}

\sphinxstepscope


\chapter{Boot Loader Entries}
\label{\detokenize{Loader-Entries:boot-loader-entries}}\label{\detokenize{Loader-Entries::doc}}
\sphinxAtStartPar
There are two kinds of systemd boot loader entries that show up in the boot menu.
\begin{itemize}
\item {} 
\sphinxAtStartPar
\sphinxstylestrong{type \#1}
\begin{itemize}
\item {} 
\sphinxAtStartPar
Uses a separate \sphinxstyleemphasis{loader entry file}

\item {} 
\sphinxAtStartPar
Used when there are separate kernel and initrd images or for efi tools.

\end{itemize}

\item {} 
\sphinxAtStartPar
\sphinxstylestrong{type \#2}
\begin{itemize}
\item {} 
\sphinxAtStartPar
Used with UKI images only.

\sphinxAtStartPar
UKI images have the kernel and initrd combined into a single file where
the UKI file itself is sufficient to provide the boot loader menu without
any additional loader entry file.

\end{itemize}

\end{itemize}

\sphinxAtStartPar
Since sd\sphinxhyphen{}boot uses \sphinxstyleemphasis{kernel\sphinxhyphen{}install} to do the real work, type \#1 boot loader entries are
automatically created with \sphinxstyleemphasis{bls} layout. It creates one new entry for each kernel version.

\sphinxAtStartPar
After using sd\sphinxhyphen{}boot the very first time (only), you should remove any stale (fixed) loader
entries from \sphinxstyleemphasis{/boot/loader/entries} or \sphinxstyleemphasis{/efi/loader/entries} (if that’s where they are) for
any kernel package now managed by sd\sphinxhyphen{}boot.

\sphinxAtStartPar
sd\sphinxhyphen{}boot also provides a kernel\sphinxhyphen{}install plugin that modifies the raw type\#1 loader entries.
These are generated for efi tools and for kernels using \sphinxstyleemphasis{bls} layout.

\sphinxAtStartPar
The primary difference is the \sphinxstyleemphasis{Title} which is modified to be the package name.
The title is shown in the boot menu.

\sphinxAtStartPar
By default the title is \sphinxstyleemphasis{Arch Linux} which comes from \sphinxstyleemphasis{/etc/os\sphinxhyphen{}release} for every
kernel and efi tool.

\sphinxAtStartPar
efi\sphinxhyphen{}tools additionally remove the kernel boot command line options from the entry
(since they are irrelevant) and change:

\begin{sphinxVerbatim}[commandchars=\\\{\}]
\PYG{n}{linux} \PYG{o}{\PYGZlt{}}\PYG{n}{tool}\PYG{o}{\PYGZgt{}}\PYG{o}{.}\PYG{n}{efi}
\end{sphinxVerbatim}

\sphinxAtStartPar
to:

\begin{sphinxVerbatim}[commandchars=\\\{\}]
\PYG{n}{efi} \PYG{o}{\PYGZlt{}}\PYG{n}{tool}\PYG{o}{\PYGZgt{}}\PYG{o}{.}\PYG{n}{efi}
\end{sphinxVerbatim}

\sphinxAtStartPar
It also means that the systemd\sphinxhyphen{}boot \sphinxstyleemphasis{loader.conf} file located in the \sphinxstyleemphasis{EFI}
needs to be updated as well when the default kernel to boot is one of those managed by sd\sphinxhyphen{}boot.

\sphinxAtStartPar
Lets show how we can use this so that a \sphinxstyleemphasis{linux\sphinxhyphen{}custom} kernel is booted by default.

\sphinxAtStartPar
The loader entries are located in:

\begin{sphinxVerbatim}[commandchars=\\\{\}]
\PYG{o}{\PYGZlt{}}\PYG{n}{EFI}\PYG{o}{\PYGZgt{}}\PYG{o}{/}\PYG{n}{loader}\PYG{o}{/}\PYG{n}{entries}\PYG{o}{/}\PYG{o}{\PYGZlt{}}\PYG{n}{machine}\PYG{o}{\PYGZhy{}}\PYG{n+nb}{id}\PYG{o}{\PYGZgt{}}\PYG{o}{\PYGZhy{}}\PYG{o}{\PYGZlt{}}\PYG{n}{kernel}\PYG{o}{\PYGZhy{}}\PYG{n}{version}\PYG{o}{\PYGZgt{}}\PYG{o}{.}\PYG{n}{conf}
\end{sphinxVerbatim}

\sphinxAtStartPar
for example:

\begin{sphinxVerbatim}[commandchars=\\\{\}]
\PYG{o}{\PYGZlt{}}\PYG{n}{EFI}\PYG{o}{\PYGZgt{}}\PYG{o}{/}\PYG{n}{loader}\PYG{o}{/}\PYG{n}{entries}\PYG{o}{/}\PYG{o}{\PYGZlt{}}\PYG{n}{machine}\PYG{o}{\PYGZhy{}}\PYG{n+nb}{id}\PYG{o}{\PYGZgt{}}\PYG{o}{\PYGZhy{}}\PYG{l+m+mf}{7.2}\PYG{l+m+mf}{.0}\PYG{o}{\PYGZhy{}}\PYG{n}{custom}\PYG{o}{\PYGZhy{}}\PYG{l+m+mf}{1.}\PYG{n}{conf}
\end{sphinxVerbatim}

\sphinxAtStartPar
We would match this entry in \sphinxstyleemphasis{loader.conf} to be the default kernel using:

\begin{sphinxVerbatim}[commandchars=\\\{\}]
\PYG{c+c1}{\PYGZsh{} loader.conf}
\PYG{n}{default} \PYG{o}{*}\PYG{o}{\PYGZhy{}}\PYG{n}{custom}\PYG{o}{\PYGZhy{}}\PYG{o}{*}
\PYG{n}{timeout} \PYG{l+m+mi}{5}
\PYG{n}{editor}  \PYG{n}{yes}
\end{sphinxVerbatim}

\sphinxAtStartPar
The wildcards are standard shell file globbing used to match the loader entry filename.
AFter making this change, its helpful to verify everything is correct by running:

\begin{sphinxVerbatim}[commandchars=\\\{\}]
bootctl
bootctl\PYG{+w}{ }list
\end{sphinxVerbatim}

\sphinxAtStartPar
For \sphinxstyleemphasis{uki} layout there are no loader entry files but the comments on \sphinxstyleemphasis{loader.conf} apply
similarly. The \sphinxstyleemphasis{loader.conf} example should work fine for BLS and UKI layouts.

\sphinxstepscope


\chapter{Layout: BLS vs UKI}
\label{\detokenize{Layout:layout-bls-vs-uki}}\label{\detokenize{Layout::doc}}
\sphinxAtStartPar
First off, I prefer \sphinxstyleemphasis{uki} layout. It is simpler, does not require separate loader entry
files, and both kernel and initrd are signed since they are in one file. The \sphinxstyleemphasis{uki} file
itself is an \sphinxstyleemphasis{efi} file and can therefore be directly booted without the need of a boot manager
should that ever be needed.

\sphinxAtStartPar
For this reason the dafault \sphinxstyleemphasis{install.conf} file provided by \sphinxstyleemphasis{sd\sphinxhyphen{}boot} uses \sphinxstyleemphasis{uki} layout.

\sphinxAtStartPar
UKI layout uses:

\begin{sphinxVerbatim}[commandchars=\\\{\}]
\PYGZdl{}BOOT/EFI/Linux/\PYGZlt{}machine\PYGZhy{}id\PYGZgt{}\PYGZhy{}\PYGZlt{}kernel\PYGZhy{}version\PYGZgt{}.efi
\end{sphinxVerbatim}

\sphinxAtStartPar
Of particular note is that there are no loader entry files in UKI layout.

\sphinxAtStartPar
BLS layout uses:

\begin{sphinxVerbatim}[commandchars=\\\{\}]
\PYGZdl{}BOOT/\PYGZlt{}machine\PYGZhy{}id\PYGZgt{}/\PYGZlt{}kernel\PYGZhy{}version\PYGZgt{}/linux
\PYGZdl{}BOOT/\PYGZlt{}machine\PYGZhy{}id\PYGZgt{}/\PYGZlt{}kernel\PYGZhy{}version\PYGZgt{}/initrd
\PYGZdl{}BOOT/loader/entries/\PYGZlt{}\PYGZlt{}machine\PYGZhy{}id\PYGZgt{}\PYGZhy{}\PYGZlt{}kernel\PYGZhy{}version\PYGZgt{}.conf
\end{sphinxVerbatim}

\sphinxAtStartPar
kernel install uses settings from \sphinxstyleemphasis{/usr/lib/kernel} which may be over\sphinxhyphen{}ridden
with files in \sphinxstyleemphasis{/etc/kernel}.

\sphinxAtStartPar
sd\sphinxhyphen{}boot installs:

\begin{sphinxVerbatim}[commandchars=\\\{\}]
\PYG{o}{/}\PYG{n}{usr}\PYG{o}{/}\PYG{n}{lib}\PYG{o}{/}\PYG{n}{kernel}\PYG{o}{/}\PYG{n}{install}\PYG{o}{.}\PYG{n}{conf}\PYG{o}{.}\PYG{n}{d}\PYG{o}{/}\PYG{l+m+mi}{010}\PYG{o}{\PYGZhy{}}\PYG{n}{sd}\PYG{o}{\PYGZhy{}}\PYG{n}{boot}\PYG{o}{\PYGZhy{}}\PYG{n}{install}\PYG{o}{.}\PYG{n}{conf}
\end{sphinxVerbatim}

\sphinxAtStartPar
which contains:

\begin{sphinxVerbatim}[commandchars=\\\{\}]
\PYG{n}{layout}\PYG{o}{=}\PYG{n}{uki}
\PYG{n}{initrd\PYGZus{}generator}\PYG{o}{=}\PYG{n}{dracut}
\PYG{n}{uki\PYGZus{}generator}\PYG{o}{=}\PYG{n}{ukify}
\end{sphinxVerbatim}

\sphinxAtStartPar
To switch back to BLS from UKI layout create or edit the file:

\begin{sphinxVerbatim}[commandchars=\\\{\}]
\PYG{o}{/}\PYG{n}{etc}\PYG{o}{/}\PYG{n}{kernel}\PYG{o}{/}\PYG{n}{install}\PYG{o}{.}\PYG{n}{conf}
\end{sphinxVerbatim}

\sphinxAtStartPar
with \sphinxstyleemphasis{layout=bls}.
Then install the kernel package again.

\sphinxAtStartPar
BLS layout uses \sphinxstyleemphasis{type \#1} loader entries. So if layout is changed from \sphinxstyleemphasis{bls} back to \sphinxstyleemphasis{uki}
neither the old BLS loader entry files nor the BLS kernel and initrd are automatically removed.
These will need to be manually removed.

\begin{sphinxVerbatim}[commandchars=\\\{\}]
rm \PYGZdl{}BOOT/loader/entries/\PYGZlt{}\PYGZlt{}machine\PYGZhy{}id\PYGZgt{}\PYGZhy{}\PYGZlt{}kernel\PYGZhy{}version\PYGZgt{}.conf
rm \PYGZhy{}rf \PYGZdl{}BOOT/\PYGZlt{}machine\PYGZhy{}id\PYGZgt{}/\PYGZlt{}kernel\PYGZhy{}version\PYGZgt{}
\end{sphinxVerbatim}

\sphinxAtStartPar
While dracut can generate the UKI file without using ukify, this has some limitations.
As of now, I recommend using ukify.

\sphinxAtStartPar
We would like to provide an additional option \sphinxstyleemphasis{OSRelease=} to ukify with a
modified version of os\sphinxhyphen{}release having \sphinxstyleemphasis{PRETTY\_NAME} and \sphinxstyleemphasis{BUILD\_ID}  leading to
a more user friendly boot menu title. This option is available in the ukify
command line tool, but kernel\sphinxhyphen{}install does not call this, instead it imports
the ukify python module and calls the functions directly. The kernel\sphinxhyphen{}install
code does not provide for the \sphinxstyleemphasis{OSRelease=} option.

\sphinxAtStartPar
At this time, however, there is no clean way to do this that I could find.
Hopefully kernel\sphinxhyphen{}install will allow this in the future. In the meantime
the boot menu items in \sphinxstyleemphasis{uki} mode are precise but quite long.

\sphinxstepscope


\chapter{Signing Kernels \& Secure Boot}
\label{\detokenize{Signed-Kernels:signing-kernels-secure-boot}}\label{\detokenize{Signed-Kernels::doc}}
\sphinxAtStartPar
Kernels will be signed if keys are available thanks to kernel\sphinxhyphen{}install and sbctl.
Many thanks to Foxboron for writing sbctl which simplifies things enormously.

\sphinxAtStartPar
This is a very brief overview.
Please see Arch wiki and man pages for details on secure boot and kernel signing.


\section{Signing Kernel}
\label{\detokenize{Signed-Kernels:signing-kernel}}
\sphinxAtStartPar
The sbctl package must be installed.

\sphinxAtStartPar
To have the kernels signed by local keys, first create the keys:

\begin{sphinxVerbatim}[commandchars=\\\{\}]
sbctl\PYG{+w}{ }create\PYGZhy{}keys
\end{sphinxVerbatim}

\sphinxAtStartPar
You may be prompted to \sphinxstyleemphasis{migrate} an older install, if so then do it:

\begin{sphinxVerbatim}[commandchars=\\\{\}]
sbctl\PYG{+w}{ }setup\PYG{+w}{ }\PYGZhy{}\PYGZhy{}migrate
\end{sphinxVerbatim}

\sphinxAtStartPar
Thats all that’s required. This will sign the kernel image. Note that if the
layout is \sphinxstyleemphasis{bls} then the initrd is separate and unsigned.

\sphinxAtStartPar
You may want to change to UKI where the initrd and the kernel image are packaged into
a single file which gets signed. To change the layout simply edit
\sphinxstyleemphasis{/etc/kernel/install.conf} and change the layout to uki:

\begin{sphinxVerbatim}[commandchars=\\\{\}]
\PYG{n+nv}{layout}\PYG{o}{=}uki
\PYG{n+nv}{initrd\PYGZus{}generator}\PYG{o}{=}dracut
\PYG{n+nv}{uki\PYGZus{}generator}\PYG{o}{=}ukify
\end{sphinxVerbatim}

\sphinxAtStartPar
That’s all that’s required to have a unified kernel image signed by the local keys.


\section{Signing BLS vs UKI}
\label{\detokenize{Signed-Kernels:signing-bls-vs-uki}}
\sphinxAtStartPar
When signing a kernel, the preferred way is to use uki layout as the initrd is part of
the uki image that gets signed. In bls layout, the kernel is signed but the initrd
is not.

\sphinxAtStartPar
When using bls layout, the kernel signing is handled by the sbct plugin, \sphinxstyleemphasis{/usr/lib/91\sphinxhyphen{}sbctl.install}.
When using uku layout, signing of the uki image is handled by ukify.

\sphinxAtStartPar
Best I can tell, the sbctl plugin does not to attempt to sign the image a second time. However,
being cautious, sd\sphinxhyphen{}boot will de\sphinxhyphen{}activate the sbctl plugin in this case to eliminate any possibility
of a second signing happening.


\section{Secure Boot}
\label{\detokenize{Signed-Kernels:secure-boot}}
\sphinxAtStartPar
\sphinxstylestrong{Big caveat}. Per the wiki, using local keys can lead to significant problems and may even break
the machine.

\sphinxAtStartPar
To use secure boot, the keys must be enrolled and secure boot activated in the UEFI Bios.

\sphinxAtStartPar
To enroll the signing keys the bios needs to have secure boot set to \sphinxstyleemphasis{setup} mode and then
boot up the machine and use \sphinxstyleemphasis{sbctl} to enroll the keys.

\begin{sphinxVerbatim}[commandchars=\\\{\}]
sbctl\PYG{+w}{ }enroll\PYGZhy{}keys\PYG{+w}{ }\PYGZhy{}\PYGZhy{}microsoft
\end{sphinxVerbatim}

\sphinxAtStartPar
At this point the machine is in secure boot mode and will only boot signed kernels.

\sphinxstepscope


\chapter{Dracut options}
\label{\detokenize{Dracut-Options:dracut-options}}\label{\detokenize{Dracut-Options::doc}}
\sphinxAtStartPar
sd\sphinxhyphen{}boot provides a default dracut config file:

\begin{sphinxVerbatim}[commandchars=\\\{\}]
/usr/lib/dracut/dracut.conf.d/010\PYGZhy{}sd\PYGZhy{}boot\PYGZhy{}dracut.conf
\end{sphinxVerbatim}

\sphinxAtStartPar
If changes are needed, then over\sphinxhyphen{}rides can be provided as usual using a drop in file:

\begin{sphinxVerbatim}[commandchars=\\\{\}]
\PYG{o}{/}\PYG{n}{etc}\PYG{o}{/}\PYG{n}{dracut}\PYG{o}{.}\PYG{n}{conf}\PYG{o}{.}\PYG{n}{d}\PYG{o}{/}\PYG{n}{xxx}\PYG{o}{.}\PYG{n}{conf}
\end{sphinxVerbatim}

\sphinxAtStartPar
since entries found in \sphinxstyleemphasis{/etc}/ take precedence over \sphinxstyleemphasis{/usr/lib/} as usual.
Files within each directory are processed in alphanumeric order. For example a file called 020\sphinxhyphen{}dracut.conf
will over\sphinxhyphen{}ride any settings 010\sphinxhyphen{}dracut.conf in the same directory.
The last file read by dracut sets the options used.

\sphinxAtStartPar
This treatment of configuration file precedence is standard on linux, and kernel\sphinxhyphen{}install follows the
same rules for it’s configuration files.

\sphinxstepscope


\chapter{Efi Filesystem Drivers}
\label{\detokenize{Efi-Filesystem-Drivers:efi-filesystem-drivers}}\label{\detokenize{Efi-Filesystem-Drivers::doc}}
\sphinxAtStartPar
While the \sphinxstyleemphasis{EFI} is always Fat\sphinxhyphen{}32, a \sphinxstyleemphasis{/boot} partition may use other filesystems.
In this case the boot loader requires efi filesystem drivers that provided by the \sphinxstyleemphasis{efifs} package.

\sphinxAtStartPar
When this package is installed the driver files are located in the directory \sphinxstyleemphasis{/usr/lib/efifs\sphinxhyphen{}x86/},
sd\sphinxhyphen{}boot detects this, via an alpm hook, and installs them into:

\begin{sphinxVerbatim}[commandchars=\\\{\}]
\PYG{o}{\PYGZlt{}}\PYG{n}{EFI}\PYG{o}{\PYGZgt{}}\PYG{o}{/}\PYG{n}{EFI}\PYG{o}{/}\PYG{n}{systemd}\PYG{o}{/}\PYG{n}{drivers}\PYG{o}{/}
\end{sphinxVerbatim}

\sphinxAtStartPar
where systemd\sphinxhyphen{}boot expects to find them.

\sphinxAtStartPar
These can also be manually installed (or removed) from the command line using:

\begin{sphinxVerbatim}[commandchars=\\\{\}]
\PYG{o}{/}\PYG{n}{usr}\PYG{o}{/}\PYG{n}{lib}\PYG{o}{/}\PYG{n}{sd}\PYG{o}{\PYGZhy{}}\PYG{n}{boot}\PYG{o}{/}\PYG{n}{sd}\PYG{o}{\PYGZhy{}}\PYG{n}{boot}\PYG{o}{\PYGZhy{}}\PYG{n}{efifs}\PYG{o}{\PYGZhy{}}\PYG{n}{update} \PYG{n}{add}
\end{sphinxVerbatim}

\sphinxAtStartPar
or:

\begin{sphinxVerbatim}[commandchars=\\\{\}]
\PYG{o}{/}\PYG{n}{usr}\PYG{o}{/}\PYG{n}{lib}\PYG{o}{/}\PYG{n}{sd}\PYG{o}{\PYGZhy{}}\PYG{n}{boot}\PYG{o}{/}\PYG{n}{sd}\PYG{o}{\PYGZhy{}}\PYG{n}{boot}\PYG{o}{\PYGZhy{}}\PYG{n}{efifs}\PYG{o}{\PYGZhy{}}\PYG{n}{update} \PYG{n}{remove}
\end{sphinxVerbatim}

\sphinxstepscope


\chapter{Skipping Kernel Plugins}
\label{\detokenize{Skipping-Kernel-Plugins:skipping-kernel-plugins}}\label{\detokenize{Skipping-Kernel-Plugins::doc}}
\sphinxAtStartPar
kernel\sphinxhyphen{}install, by default, calls every plugin.
There may be times when there’s a need to de\sphinxhyphen{}activate a plugin. This is accomplished using
a setting in sd\sphinxhyphen{}boot config file:

\begin{sphinxVerbatim}[commandchars=\\\{\}]
\PYG{o}{/}\PYG{n}{etc}\PYG{o}{/}\PYG{n}{sd}\PYG{o}{\PYGZhy{}}\PYG{n}{boot}\PYG{o}{/}\PYG{n}{config}\PYG{o}{.}\PYG{n}{yaml}
\end{sphinxVerbatim}

\sphinxAtStartPar
Set \sphinxstyleemphasis{skip\_kernel\_plugins} to  the list of plugins to be skippedi.
Each must be a full pathname:

\begin{sphinxVerbatim}[commandchars=\\\{\}]
\PYG{n}{skip\PYGZus{}kernel\PYGZus{}plugins}\PYG{p}{:}
    \PYG{o}{\PYGZhy{}} \PYG{o}{/}\PYG{n}{usr}\PYG{o}{/}\PYG{n}{lib}\PYG{o}{/}\PYG{n}{kernel}\PYG{o}{/}\PYG{n}{install}\PYG{o}{.}\PYG{n}{d}\PYG{o}{/}\PYG{l+m+mi}{91}\PYG{o}{\PYGZhy{}}\PYG{n}{xxx}\PYG{o}{.}\PYG{n}{install}
    \PYG{o}{\PYGZhy{}} \PYG{o}{/}\PYG{n}{usr}\PYG{o}{/}\PYG{n}{lib}\PYG{o}{/}\PYG{n}{kernel}\PYG{o}{/}\PYG{n}{install}\PYG{o}{.}\PYG{n}{d}\PYG{o}{/}\PYG{l+m+mi}{91}\PYG{o}{\PYGZhy{}}\PYG{n}{yyy}\PYG{o}{.}\PYG{n}{install}
\end{sphinxVerbatim}

\sphinxAtStartPar
sd\sphinxhyphen{}boot effects this by passing the plugin list to be used to kernel\sphinxhyphen{}install using the environment variable,
KERNEL\_INSTALL\_PLUGINS.

\sphinxAtStartPar
This environment variable contains the list of all plugins to be invoked by kernel\sphinxhyphen{}install.
The list excludes any of those configured as above to be skipped.

\sphinxstepscope


\chapter{Bootable efi tools}
\label{\detokenize{Efi-Tools:bootable-efi-tools}}\label{\detokenize{Efi-Tools::doc}}
\sphinxAtStartPar
Please ensure the Arch package name is listed in /etc/sd\sphinxhyphen{}boot/efi\sphinxhyphen{}tools to let sd\sphinxhyphen{}boot
know which efi tools it is permitted to install.

\sphinxAtStartPar
sd\sphinxhyphen{}boot supports bootable efi tools such as \sphinxstyleemphasis{efi\sphinxhyphen{}shell} provided by the \sphinxstyleemphasis{edk2\sphinxhyphen{}shell} package.

\sphinxAtStartPar
The alpm hook \sphinxstyleemphasis{99\sphinxhyphen{}sd\sphinxhyphen{}boot\sphinxhyphen{}efi\sphinxhyphen{}tool\sphinxhyphen{}install.hook} provides the trigger based on package name
and the location of the efi file itself is found in

\begin{sphinxVerbatim}[commandchars=\\\{\}]
/etc/sd\PYGZhy{}boot/edk2\PYGZhy{}shell.image
\end{sphinxVerbatim}

\sphinxAtStartPar
This file contains the efi file path. It may also contain comments (lines starting with \#).

\sphinxAtStartPar
Loader entries for efi tools are not kernels, and sd\sphinxhyphen{}boot uses a kernel\sphinxhyphen{}install plugin
to modify the entry appropriately.

\sphinxAtStartPar
To add any efi tool and have it \sphinxstyleemphasis{just work}, 2 things are needed.
\begin{itemize}
\item {} 
\sphinxAtStartPar
ALPM hook files to install and remove the tool from \$BOOT.

\item {} 
\sphinxAtStartPar
A file containing the path to the bootable efi file itself.

\end{itemize}

\sphinxAtStartPar
Simplest way to add a new efi tool, is to use the efi shell as templates and modify appropriately.

\sphinxAtStartPar
For example, if the efi tool package is called XXX\sphinxhyphen{}efi
then the alpm hook install file, installed (as usual) in:

\begin{sphinxVerbatim}[commandchars=\\\{\}]
\PYG{o}{/}\PYG{n}{usr}\PYG{o}{/}\PYG{n}{share}\PYG{o}{/}\PYG{n}{libalpm}\PYG{o}{/}\PYG{l+m+mi}{99}\PYG{o}{\PYGZhy{}}\PYG{n}{XXX}\PYG{o}{\PYGZhy{}}\PYG{n}{efi}\PYG{o}{\PYGZhy{}}\PYG{n}{install}\PYG{o}{.}\PYG{n}{hook}
\end{sphinxVerbatim}

\sphinxAtStartPar
should contain:

\begin{sphinxVerbatim}[commandchars=\\\{\}]
\PYG{p}{[}\PYG{n}{Trigger}\PYG{p}{]}
 \PYG{n}{Type} \PYG{o}{=} \PYG{n}{Package}
 \PYG{n}{Operation} \PYG{o}{=} \PYG{n}{Install}
 \PYG{n}{Operation} \PYG{o}{=} \PYG{n}{Upgrade}
 \PYG{n}{Target} \PYG{o}{=} \PYG{n}{XXX}\PYG{o}{\PYGZhy{}}\PYG{n}{efi}

 \PYG{p}{[}\PYG{n}{Action}\PYG{p}{]}
 \PYG{n}{Description} \PYG{o}{=} \PYG{n}{sd}\PYG{o}{\PYGZhy{}}\PYG{n}{boot}\PYG{p}{:} \PYG{n}{installing} \PYG{n}{an} \PYG{n}{efi} \PYG{n}{tool} \PYG{n}{to} \PYG{n}{EFI}
 \PYG{n}{When} \PYG{o}{=} \PYG{n}{PostTransaction}
 \PYG{n}{Exec} \PYG{o}{=} \PYG{o}{/}\PYG{n}{usr}\PYG{o}{/}\PYG{n}{lib}\PYG{o}{/}\PYG{n}{sd}\PYG{o}{\PYGZhy{}}\PYG{n}{boot}\PYG{o}{/}\PYG{n}{sd}\PYG{o}{\PYGZhy{}}\PYG{n}{boot}\PYG{o}{\PYGZhy{}}\PYG{n}{efi}\PYG{o}{\PYGZhy{}}\PYG{n}{tool}\PYG{o}{\PYGZhy{}}\PYG{n}{update} \PYG{n}{add}
 \PYG{n}{NeedsTargets}
\end{sphinxVerbatim}

\sphinxAtStartPar
The remove file:

\begin{sphinxVerbatim}[commandchars=\\\{\}]
\PYG{o}{/}\PYG{n}{usr}\PYG{o}{/}\PYG{n}{share}\PYG{o}{/}\PYG{n}{libalpm}\PYG{o}{/}\PYG{l+m+mi}{70}\PYG{o}{\PYGZhy{}}\PYG{n}{XXX}\PYG{o}{\PYGZhy{}}\PYG{n}{efi}\PYG{o}{\PYGZhy{}}\PYG{n}{remove}\PYG{o}{.}\PYG{n}{hook}
\end{sphinxVerbatim}

\sphinxAtStartPar
should contain:

\begin{sphinxVerbatim}[commandchars=\\\{\}]
\PYG{p}{[}\PYG{n}{Trigger}\PYG{p}{]}
\PYG{n}{Type} \PYG{o}{=} \PYG{n}{Package}
\PYG{n}{Operation} \PYG{o}{=} \PYG{n}{Remove}
\PYG{n}{Target} \PYG{o}{=} \PYG{n}{XXX}\PYG{o}{\PYGZhy{}}\PYG{n}{efi}

\PYG{p}{[}\PYG{n}{Action}\PYG{p}{]}
\PYG{n}{Description} \PYG{o}{=} \PYG{n}{sd}\PYG{o}{\PYGZhy{}}\PYG{n}{boot}\PYG{p}{:} \PYG{n}{Removing} \PYG{n}{efi} \PYG{n}{tool} \PYG{k+kn}{from}\PYG{+w}{ }\PYG{n+nn}{ESP}
\PYG{n}{When} \PYG{o}{=} \PYG{n}{PreTransaction}
\PYG{n}{Exec} \PYG{o}{=} \PYG{o}{/}\PYG{n}{usr}\PYG{o}{/}\PYG{n}{lib}\PYG{o}{/}\PYG{n}{sd}\PYG{o}{\PYGZhy{}}\PYG{n}{boot}\PYG{o}{/}\PYG{n}{sd}\PYG{o}{\PYGZhy{}}\PYG{n}{boot}\PYG{o}{\PYGZhy{}}\PYG{n}{efi}\PYG{o}{\PYGZhy{}}\PYG{n}{tool}\PYG{o}{\PYGZhy{}}\PYG{n}{update} \PYG{n}{remove}
\PYG{n}{NeedsTargets}
\end{sphinxVerbatim}

\sphinxAtStartPar
The last file provides the location of the efi file itself. This should be located in:

\begin{sphinxVerbatim}[commandchars=\\\{\}]
\PYG{o}{/}\PYG{n}{etc}\PYG{o}{/}\PYG{n}{sd}\PYG{o}{\PYGZhy{}}\PYG{n}{boot}\PYG{o}{/}\PYG{n}{XXX}\PYG{o}{\PYGZhy{}}\PYG{n}{efi}\PYG{o}{.}\PYG{n}{image}
\end{sphinxVerbatim}

\sphinxAtStartPar
and contain the path itself:

\begin{sphinxVerbatim}[commandchars=\\\{\}]
\PYG{o}{/}\PYG{n}{usr}\PYG{o}{/}\PYG{n}{share}\PYG{o}{/}\PYG{n}{XXX}\PYG{o}{\PYGZhy{}}\PYG{n}{efi}\PYG{o}{/}\PYG{n}{x64}\PYG{o}{/}\PYG{n}{xxx}\PYG{o}{\PYGZhy{}}\PYG{n}{tool}\PYG{o}{.}\PYG{n}{efi}
\end{sphinxVerbatim}


\section{A note on memtest86 plus}
\label{\detokenize{Efi-Tools:a-note-on-memtest86-plus}}
\sphinxAtStartPar
If sd\sphinxhyphen{}boot is used to install this please do not use the (current)
version of \sphinxstyleemphasis{memtest86+\sphinxhyphen{}efi} in the Arch repo. As of this writing, this package installs
files directly into \sphinxstyleemphasis{/boot} which violates the basic requierments.
It is also somewhat out of date.

\sphinxAtStartPar
I may provide a compatible version to the AUR at some point if its helpful.
Note that the latest open source version from memtest.org is 8.x. There is
also a non\sphinxhyphen{}open source (commercial) version) 11.x in the AUR that installs to /usr/share.
The latest open source version is v8.10 at this time.

\sphinxstepscope


\chapter{Development}
\label{\detokenize{Development:development}}\label{\detokenize{Development::doc}}

\section{C\sphinxhyphen{}Code Language Choice}
\label{\detokenize{Development:c-code-language-choice}}
\sphinxAtStartPar
While I was tempted to do this in python, partly because systemd chose to use it for
\sphinxstyleemphasis{/usr/lib/kernel/install.d/60\sphinxhyphen{}ukify.install}, I decided to use C.

\sphinxAtStartPar
It’s important to be able to do initial testing and validation as non\sphinxhyphen{}root user
and of course without touching any important files in any actual root directory.

\sphinxAtStartPar
While the the C\sphinxhyphen{}code version is a bit more work and a little more complicated to create,
the benefit, in my view, is that the c\sphinxhyphen{}code is much easier to test and debug
and keep organized than bash version.

\sphinxAtStartPar
C also has the advantage that the compiled code has almost no dependencies and
is highly performant. Way faster than bash or python. For a system utility this
is a pretty nice benefit.


\section{Testing C\sphinxhyphen{}Code}
\label{\detokenize{Development:testing-c-code}}
\sphinxAtStartPar
Testing and development are done using a \sphinxstyleemphasis{Testing} directory that is writable by non\sphinxhyphen{}root user.
This is where kernels will be installed, loader entries updated and so on. The tools
are installed into this tree and are statically linked.

\sphinxAtStartPar
Otherwise the environment is kept clean other than the \sphinxstyleemphasis{PATH} variable since dracut requires it.

\sphinxAtStartPar
This allows all programs to be tested by leveraging kernel\sphinxhyphen{}install’s ability to work
completely in the testing root tree.

\sphinxAtStartPar
The environment variable \sphinxstyleemphasis{SDB\_DEV\_TEST} activates the test usage, allowing
testing and debugging to be done completely within the non\sphinxhyphen{}root test tree.

\sphinxAtStartPar
There are two test sets provided under \sphinxstyleemphasis{src/c\sphinxhyphen{}code/scripts}. Both test sets
should be run from the \sphinxstyleemphasis{src} directory.
\begin{itemize}
\item {} 
\sphinxAtStartPar
./tests/scripts/run\sphinxhyphen{}test\sphinxhyphen{}suite

\sphinxAtStartPar
This runs the the tools with \sphinxstyleemphasis{root} set to \sphinxstyleemphasis{Testing/\_\_root\_\_}.
This is run as part of the build check process.
It also runs the tools under valgrind.

\item {} 
\sphinxAtStartPar
./tests/scripts/static\sphinxhyphen{}analysis:

\sphinxAtStartPar
It runs static code analysis using cppcheck and clang\sphinxhyphen{}tidy.
Note that kernel\sphinxhyphen{}install always runs \sphinxstyleemphasis{chown} and since all tests are run as ordinary
(non\sphinxhyphen{}root) user this leads to some warning messages landing in the testing logs.
These are benign as kernel\sphinxhyphen{}install ignores them. They dont happen in production
where kernel\sphinxhyphen{}install is running as root and chown is permitted.

\sphinxAtStartPar
For each of these tests, the logs save stdout, stderr and the exit status of the tool
along with the output from valgrind.

\end{itemize}

\sphinxAtStartPar
The image and initrd files will be installed in:

\begin{sphinxVerbatim}[commandchars=\\\{\}]
\PYG{n}{Testing}\PYG{o}{/}\PYG{n}{\PYGZus{}\PYGZus{}root\PYGZus{}\PYGZus{}}\PYG{o}{/}\PYG{n}{boot}\PYG{o}{/}\PYG{o}{\PYGZlt{}}\PYG{n}{machine}\PYG{o}{\PYGZhy{}}\PYG{n+nb}{id}\PYG{o}{\PYGZgt{}}\PYG{o}{/}\PYG{o}{\PYGZlt{}}\PYG{n}{kernel}\PYG{o}{\PYGZhy{}}\PYG{n}{version}\PYG{o}{\PYGZgt{}}\PYG{o}{/}
\end{sphinxVerbatim}

\sphinxAtStartPar
while the loader entry files, for efi tools and kernels in \sphinxstyleemphasis{bls} layout will be in:

\begin{sphinxVerbatim}[commandchars=\\\{\}]
\PYG{n}{Testing}\PYG{o}{/}\PYG{n}{\PYGZus{}\PYGZus{}root\PYGZus{}\PYGZus{}}\PYG{o}{/}\PYG{n}{boot}\PYG{o}{/}\PYG{n}{loader}\PYG{o}{/}\PYG{n}{entries}\PYG{o}{/}\PYG{o}{\PYGZlt{}}\PYG{n}{machine}\PYG{o}{\PYGZhy{}}\PYG{n+nb}{id}\PYG{o}{\PYGZgt{}}\PYG{o}{\PYGZhy{}}\PYG{o}{\PYGZlt{}}\PYG{n}{kernel}\PYG{o}{\PYGZhy{}}\PYG{n}{version}\PYG{o}{\PYGZgt{}}\PYG{o}{.}\PYG{n}{conf}
\end{sphinxVerbatim}

\sphinxAtStartPar
You may notice that the loader entries have a longer path to the kernel image and initrd.

\sphinxAtStartPar
This is normal.

\sphinxAtStartPar
When run in a test tree, kernel\sphinxhyphen{}install identifies the \sphinxstyleemphasis{mount} point and removes it from
the front of the path. In test mode where the test directory is not actually a mount point
(such as /boot) the pathname written to the loader entry will include more elements.
This is fine and when run in producion the image file will simply be \sphinxstyleemphasis{linux} instead of
\sphinxstyleemphasis{/long/path/to/linux}. The same is true for the initrd file as well.


\section{Development and Debugging}
\label{\detokenize{Development:development-and-debugging}}
\sphinxAtStartPar
Debuggable executables can be installed into the test root directory.
From the \sphinxstyleemphasis{src} directory:

\begin{sphinxVerbatim}[commandchars=\\\{\}]
./do\PYGZhy{}dev\PYGZhy{}build
\PYG{n+nb}{cd}\PYG{+w}{ }tests/Testing
\PYG{n+nb}{export}\PYG{+w}{ }\PYG{n+nv}{SDB\PYGZus{}DEV\PYGZus{}TEST}\PYG{o}{=}\PYG{n+nb}{true}
\end{sphinxVerbatim}

\sphinxAtStartPar
Then executables can now be debugged using, for example:

\begin{sphinxVerbatim}[commandchars=\\\{\}]
gdb\PYG{+w}{ }./\PYGZus{}\PYGZus{}root\PYGZus{}\PYGZus{}/usr/bin/sd\PYGZhy{}boot\PYGZhy{}kernel\PYGZhy{}update
\end{sphinxVerbatim}

\sphinxAtStartPar
And for this executable for example, within gdb use:

\begin{sphinxVerbatim}[commandchars=\\\{\}]
\PYG{n}{run} \PYG{n}{add} \PYG{n}{linux}
\end{sphinxVerbatim}

\sphinxAtStartPar
Plugins are called by kernel\sphinxhyphen{}install but may be run manually as well:
\begin{itemize}
\item {} 
\sphinxAtStartPar
./plugin\sphinxhyphen{}tests/run\sphinxhyphen{}loaderentry\sphinxhyphen{}efi (for bls plugin\sphinxhyphen{}tests/run\sphinxhyphen{}loaderentry\sphinxhyphen{}kernel)

\sphinxAtStartPar
Standalone tests of the plugin that modify the raw loader entry files.
The scripts run the tests under valgrind, but there is an option
to run in the debugger as well.

\end{itemize}

\sphinxstepscope


\chapter{Appendix}
\label{\detokenize{Appendix:appendix}}\label{\detokenize{Appendix::doc}}

\section{Potential Todo Items}
\label{\detokenize{Appendix:potential-todo-items}}\begin{itemize}
\item {} 
\sphinxAtStartPar
When switching layout from bls to uki, previous kernel is not removed
since it was installed using bls layout while \sphinxstyleemphasis{kernel\sphinxhyphen{}install remove}
is now using uki layout. This leaves the previous bls kernel install instead of
removing it. Same would be true switching layout from uki back to bls.

\sphinxAtStartPar
This means, at least for now, that manual intervention is necessary after changing layout
to avoid leaving un\sphinxhyphen{}needed files. While it is benign they do take disk space.

\sphinxAtStartPar
For example after changing to uki layout, check \sphinxstyleemphasis{/boot/loader/entries} and
\sphinxstyleemphasis{/boot/<machine\sphinxhyphen{}id>/} and remove the older kernel(s) that are no longer needed.

\sphinxAtStartPar
In uki mode there are no loader entry files at all and kernels are installed
in \sphinxstyleemphasis{/boot/EFI/Linux} not in \sphinxstyleemphasis{/boot/<machine\sphinxhyphen{}id>}.

\sphinxAtStartPar
Here \sphinxstyleemphasis{/boot} means either \sphinxstyleemphasis{/boot} or \sphinxstyleemphasis{/efi} as appropriate.

\item {} 
\sphinxAtStartPar
So, it might be good (nice to have) for the code be more helpful with this.

\end{itemize}


\chapter{Man Pages}
\label{\detokenize{index:man-pages}}
\sphinxstepscope


\section{sd\sphinxhyphen{}boot\sphinxhyphen{}efifs\sphinxhyphen{}update}
\label{\detokenize{_build/man-ready/sd-boot-efifs-update.8:sd-boot-efifs-update}}\label{\detokenize{_build/man-ready/sd-boot-efifs-update.8::doc}}

\subsection{Install or remove efi filesystem drivers to ESP}
\label{\detokenize{_build/man-ready/sd-boot-efifs-update.8:install-or-remove-efi-filesystem-drivers-to-esp}}

\subsubsection{SYNOPSIS}
\label{\detokenize{_build/man-ready/sd-boot-efifs-update.8:synopsis}}
\sphinxAtStartPar
\sphinxcode{\sphinxupquote{sd\sphinxhyphen{}boot\sphinxhyphen{}efifs\sphinxhyphen{}update}} \sphinxcode{\sphinxupquote{add}} | \sphinxcode{\sphinxupquote{remove}}


\subsubsection{DESCRIPTION}
\label{\detokenize{_build/man-ready/sd-boot-efifs-update.8:description}}
\sphinxAtStartPar
Install or remove efi filesystem drivers provided by \sphinxstyleemphasis{efifs} package to the \sphinxstyleemphasis{EFI} partition.
The driver files are copied from /usr/lib/efifs\sphinxhyphen{}x64.


\subsubsection{ARGUMENTS}
\label{\detokenize{_build/man-ready/sd-boot-efifs-update.8:arguments}}
\sphinxAtStartPar
\sphinxcode{\sphinxupquote{add}}
\begin{quote}

\sphinxAtStartPar
Copies the efi drivers provided by the package and installs them
to the EFI partition.
\end{quote}

\sphinxAtStartPar
\sphinxcode{\sphinxupquote{remove}}
\begin{quote}

\sphinxAtStartPar
Remove the drivers from the EFI partition. This does not touch the source files
in \sphinxstyleemphasis{/usr/lib/efifs\sphinxhyphen{}x64} which are managed by pacman.
\end{quote}


\subsubsection{FILES}
\label{\detokenize{_build/man-ready/sd-boot-efifs-update.8:files}}
\sphinxAtStartPar
\sphinxcode{\sphinxupquote{/usr/lib/sd\sphinxhyphen{}boot/sd\sphinxhyphen{}boot\sphinxhyphen{}efifs\sphinxhyphen{}update}}
\begin{quote}

\sphinxAtStartPar
Location of the executable.
\end{quote}

\sphinxAtStartPar
\sphinxcode{\sphinxupquote{<EFI>/EFI/systemd/drivers/}}
\begin{quote}

\sphinxAtStartPar
The efi partition where the drivers are installed. \sphinxstyleemphasis{<EFI>} is the mount point of the
ESP partition.
\end{quote}


\subsubsection{EXAMPLES}
\label{\detokenize{_build/man-ready/sd-boot-efifs-update.8:examples}}
\sphinxAtStartPar
To install the drivers:

\begin{sphinxVerbatim}[commandchars=\\\{\}]
\PYG{c+c1}{\PYGZsh{} /usr/lib/sd\PYGZhy{}boot/sd\PYGZhy{}boot\PYGZhy{}efifs\PYGZhy{}update add}
\end{sphinxVerbatim}

\sphinxAtStartPar
To remove the drivers:

\begin{sphinxVerbatim}[commandchars=\\\{\}]
\PYG{c+c1}{\PYGZsh{} /usr/lib/sd\PYGZhy{}boot/sd\PYGZhy{}boot\PYGZhy{}efifs\PYGZhy{}update remove}
\end{sphinxVerbatim}


\subsubsection{SEE ALSO}
\label{\detokenize{_build/man-ready/sd-boot-efifs-update.8:see-also}}
\sphinxAtStartPar
\sphinxstylestrong{sd\sphinxhyphen{}boot\sphinxhyphen{}efi\sphinxhyphen{}tool\sphinxhyphen{}update}, \sphinxstylestrong{sd\sphinxhyphen{}boot\sphinxhyphen{}kernel\sphinxhyphen{}update}, \sphinxstylestrong{sd\sphinxhyphen{}boot\sphinxhyphen{}find\sphinxhyphen{}boot\sphinxhyphen{}mounts},
\sphinxstylestrong{kernel\sphinxhyphen{}install}

\sphinxstepscope


\section{sd\sphinxhyphen{}boot\sphinxhyphen{}efi\sphinxhyphen{}tool\sphinxhyphen{}update}
\label{\detokenize{_build/man-ready/sd-boot-efi-tool-update.8:sd-boot-efi-tool-update}}\label{\detokenize{_build/man-ready/sd-boot-efi-tool-update.8::doc}}

\subsection{Install or remove an efi tool to \$BOOT}
\label{\detokenize{_build/man-ready/sd-boot-efi-tool-update.8:install-or-remove-an-efi-tool-to-boot}}

\subsubsection{SYNOPSIS}
\label{\detokenize{_build/man-ready/sd-boot-efi-tool-update.8:synopsis}}
\sphinxAtStartPar
\sphinxcode{\sphinxupquote{sd\sphinxhyphen{}boot\sphinxhyphen{}efi\sphinxhyphen{}tool\sphinxhyphen{}update}} \sphinxcode{\sphinxupquote{add}} | \sphinxcode{\sphinxupquote{remove}} | \sphinxcode{\sphinxupquote{inspect}} \sphinxcode{\sphinxupquote{<package\_name>}}


\subsubsection{DESCRIPTION}
\label{\detokenize{_build/man-ready/sd-boot-efi-tool-update.8:description}}
\sphinxAtStartPar
Install, remove, or inspect a bootable efi tool in the \$BOOT partition using kernel\sphinxhyphen{}install.
The \$BOOT partition, following the notation used by kernel\sphinxhyphen{}install, is usually one of \sphinxstyleemphasis{/boot} or \sphinxstyleemphasis{/efi}.

\sphinxAtStartPar
The package must be managed by sd\sphinxhyphen{}boot with the package name one of those listed in:

\begin{sphinxVerbatim}[commandchars=\\\{\}]
\PYG{o}{/}\PYG{n}{etc}\PYG{o}{/}\PYG{n}{sd}\PYG{o}{\PYGZhy{}}\PYG{n}{boot}\PYG{o}{/}\PYG{n}{efi}\PYG{o}{\PYGZhy{}}\PYG{n}{tool}\PYG{o}{.}\PYG{n}{packages}
\end{sphinxVerbatim}

\sphinxAtStartPar
It takes two arguments. The first argument is the operation:

\begin{sphinxVerbatim}[commandchars=\\\{\}]
\PYG{o}{*}\PYG{n}{add}\PYG{o}{*} \PYG{o}{|} \PYG{o}{*}\PYG{n}{remove}\PYG{o}{*} \PYG{o}{|} \PYG{o}{*}\PYG{n}{inspect}\PYG{o}{*}
\end{sphinxVerbatim}

\sphinxAtStartPar
followed by the package name that provides the tool.

\sphinxAtStartPar
For installs using \sphinxstyleemphasis{add}, the path of the source \sphinxstylestrong{efi} image must be provided in:

\begin{sphinxVerbatim}[commandchars=\\\{\}]
\PYG{o}{/}\PYG{n}{etc}\PYG{o}{/}\PYG{n}{sd}\PYG{o}{\PYGZhy{}}\PYG{n}{boot}\PYG{o}{/}\PYG{o}{\PYGZlt{}}\PYG{n}{package}\PYG{o}{\PYGZhy{}}\PYG{n}{name}\PYG{o}{\PYGZgt{}}\PYG{o}{.}\PYG{n}{image}
\end{sphinxVerbatim}

\sphinxAtStartPar
This file may contain shell style comments as well as the image path.

\sphinxAtStartPar
Since these tools are not kernels, they are installed using \sphinxstyleemphasis{BLS} layout, even
when the kernel layout is set to \sphinxstyleemphasis{UKI} in \sphinxstyleemphasis{/etc/kernel/install.conf}. This ensures that
the efi tools are installed to:

\begin{sphinxVerbatim}[commandchars=\\\{\}]
\PYG{o}{/}\PYG{n}{boot}\PYG{o}{/}\PYG{o}{\PYGZlt{}}\PYG{n}{machine}\PYG{o}{\PYGZhy{}}\PYG{n+nb}{id}\PYG{o}{\PYGZgt{}}\PYG{o}{/}\PYG{o}{\PYGZlt{}}\PYG{n}{package\PYGZus{}name}\PYG{o}{\PYGZgt{}}\PYG{o}{/}
\end{sphinxVerbatim}

\sphinxAtStartPar
and that boot loader entry files are appropriately created in:

\begin{sphinxVerbatim}[commandchars=\\\{\}]
\PYG{o}{/}\PYG{n}{boot}\PYG{o}{/}\PYG{n}{loader}\PYG{o}{/}\PYG{n}{entries}\PYG{o}{/}\PYG{o}{\PYGZlt{}}\PYG{n}{machine\PYGZus{}name}\PYG{o}{\PYGZhy{}}\PYG{n}{pacmage\PYGZus{}name}\PYG{o}{\PYGZgt{}}\PYG{o}{.}\PYG{n}{conf}
\end{sphinxVerbatim}


\subsubsection{ARGUMENTS}
\label{\detokenize{_build/man-ready/sd-boot-efi-tool-update.8:arguments}}
\sphinxAtStartPar
\sphinxstylestrong{Operation}:
\begin{itemize}
\item {} 
\sphinxAtStartPar
add
\begin{quote}

\sphinxAtStartPar
Install the efi tool to \$BOOT partition
\end{quote}

\item {} 
\sphinxAtStartPar
remove
\begin{quote}

\sphinxAtStartPar
Remove the drivers from \$BOOT partition.
\end{quote}

\item {} 
\sphinxAtStartPar
inspect
\begin{quote}

\sphinxAtStartPar
Disply information about an installed package.
\end{quote}

\end{itemize}

\sphinxAtStartPar
\sphinxstylestrong{Package Name}

\sphinxAtStartPar
Can be name of a package or the string “–all–” which means all efi tools managed by sd\sphinxhyphen{}boot.


\subsubsection{FILES}
\label{\detokenize{_build/man-ready/sd-boot-efi-tool-update.8:files}}
\sphinxAtStartPar
\sphinxcode{\sphinxupquote{/etc/sd\sphinxhyphen{}boot/<package\sphinxhyphen{}name>.packages}}
\begin{quote}

\sphinxAtStartPar
The package name must be listed here to permit sd\sphinxhyphen{}boot to manage the tool.
\end{quote}

\sphinxAtStartPar
\sphinxcode{\sphinxupquote{/etc/sd\sphinxhyphen{}boot/<package\sphinxhyphen{}name>.image}}
\begin{quote}

\sphinxAtStartPar
This file contains the full path of the efi tool provided by the package. This efi image
file will be copied to \$BOOT.
\end{quote}

\sphinxAtStartPar
\sphinxcode{\sphinxupquote{\$BOOT/<machine\sphinxhyphen{}id>/<package\_name>\sphinxhyphen{}<package\_version>}}
\begin{quote}

\sphinxAtStartPar
This is the path where the bootable efi tool is installed.
\end{quote}

\sphinxAtStartPar
\sphinxcode{\sphinxupquote{\$BOOT/loader/entries/<machine\sphinxhyphen{}id>\sphinxhyphen{}<package\_name>\sphinxhyphen{}<package\_version>.conf}}
\begin{quote}

\sphinxAtStartPar
This is the type \#1 boot loader entry file. It has the title that is used in boot menu
along with the path to the program to be booted.
\end{quote}


\subsubsection{EXAMPLES}
\label{\detokenize{_build/man-ready/sd-boot-efi-tool-update.8:examples}}
\sphinxAtStartPar
To install the efi shell from \sphinxstyleemphasis{edk2\sphinxhyphen{}shell} package:

\begin{sphinxVerbatim}[commandchars=\\\{\}]
\PYG{c+c1}{\PYGZsh{} sd\PYGZhy{}boot\PYGZhy{}efi\PYGZhy{}tool\PYGZhy{}update add edk2\PYGZhy{}shell}
\end{sphinxVerbatim}

\sphinxAtStartPar
where edk2\sphinxhyphen{}shell is listed in \sphinxstyleemphasis{/etc/sd\sphinxhyphen{}boot/edk2\sphinxhyphen{}shell.packages} and the pathname
is provided in the file \sphinxstyleemphasis{/etc/sd\sphinxhyphen{}boot/edk2\sphinxhyphen{}shell.image} which contains:

\begin{sphinxVerbatim}[commandchars=\\\{\}]
\PYG{o}{/}\PYG{n}{usr}\PYG{o}{/}\PYG{n}{share}\PYG{o}{/}\PYG{n}{edk2}\PYG{o}{\PYGZhy{}}\PYG{n}{shell}\PYG{o}{/}\PYG{n}{x64}\PYG{o}{/}\PYG{n}{Shell\PYGZus{}Full}\PYG{o}{.}\PYG{n}{efi}
\end{sphinxVerbatim}

\sphinxAtStartPar
The image itself is provided by the \sphinxstyleemphasis{edk2\sphinxhyphen{}shell} package.

\sphinxAtStartPar
To update all efi tools currently managed by sd\sphinxhyphen{}boot:

\begin{sphinxVerbatim}[commandchars=\\\{\}]
\PYG{n}{sd}\PYG{o}{\PYGZhy{}}\PYG{n}{boot}\PYG{o}{\PYGZhy{}}\PYG{n}{efi}\PYG{o}{\PYGZhy{}}\PYG{n}{tool}\PYG{o}{\PYGZhy{}}\PYG{n}{update} \PYG{n}{add} \PYG{l+s+s2}{\PYGZdq{}}\PYG{l+s+s2}{\PYGZhy{}\PYGZhy{}all\PYGZhy{}\PYGZhy{}}\PYG{l+s+s2}{\PYGZdq{}}
\end{sphinxVerbatim}

\sphinxstepscope


\section{sd\sphinxhyphen{}boot\sphinxhyphen{}find\sphinxhyphen{}boot\sphinxhyphen{}mounts}
\label{\detokenize{_build/man-ready/sd-boot-find-boot-mounts.8:sd-boot-find-boot-mounts}}\label{\detokenize{_build/man-ready/sd-boot-find-boot-mounts.8::doc}}

\subsection{Display ESP and XBOOTLDR partition mount points}
\label{\detokenize{_build/man-ready/sd-boot-find-boot-mounts.8:display-esp-and-xbootldr-partition-mount-points}}

\subsubsection{SYNOPSIS}
\label{\detokenize{_build/man-ready/sd-boot-find-boot-mounts.8:synopsis}}
\sphinxAtStartPar
\sphinxcode{\sphinxupquote{sd\sphinxhyphen{}boot\sphinxhyphen{}find\sphinxhyphen{}boot\sphinxhyphen{}mounts}}


\subsubsection{DESCRIPTION}
\label{\detokenize{_build/man-ready/sd-boot-find-boot-mounts.8:description}}
\sphinxAtStartPar
Displays the list of all ESP and XBOOTLDR partitions. Each is also marked
indicating any that are \sphinxstyleemphasis{current} and the \sphinxstyleemphasis{active} ESP used to boot the machine.

\sphinxAtStartPar
Each may have
a mark indicating whether that partition is the current one. The ESP
used by the motherboard to boot the machine is also saved to nvram.
This partition, the one active at boot time, is also marked.

\sphinxAtStartPar
When machine boots the boot loader on the motherboard boots from an ESP and it saves
that \sphinxstyleemphasis{active} ESP partition into nvram. The \sphinxstyleemphasis{current} or expected ESP
is the one that is designated by the boot order.

\sphinxAtStartPar
Under normal circumstances and almost always,  \sphinxstyleemphasis{current} and \sphinxstyleemphasis{active} are the same.
See ‘Quirky Esp’ section that discusses the situation when these differ and
a possible resolution.


\subsubsection{ARGUMENTS}
\label{\detokenize{_build/man-ready/sd-boot-find-boot-mounts.8:arguments}}
\sphinxAtStartPar
\sphinxcode{\sphinxupquote{none}}


\subsubsection{EXAMPLES}
\label{\detokenize{_build/man-ready/sd-boot-find-boot-mounts.8:examples}}\begin{itemize}
\item {} 
\sphinxAtStartPar
Example 1

\sphinxAtStartPar
One ESP:

\begin{sphinxVerbatim}[commandchars=\\\{\}]
                                              (c)urrent (a)ctive
        Device Mount                          (c a) type
/dev/nvme0n1p1 /boot                          (\(\pmb{\checkmark}\) \(\pmb{\checkmark}\)) EFI
\end{sphinxVerbatim}

\item {} 
\sphinxAtStartPar
Example 2:

\sphinxAtStartPar
One ESP and one XBOOTLDR:

\begin{sphinxVerbatim}[commandchars=\\\{\}]
   Device Mount                          (c a) type
/dev/sda1 /efi                           (\(\pmb{\checkmark}\) \(\pmb{\checkmark}\)) EFI
/dev/sda5 /boot                          (\(\pmb{\checkmark}\) \(\pmb{\checkmark}\)) XBOOTLDR
\end{sphinxVerbatim}

\end{itemize}
\begin{quote}

\sphinxAtStartPar
and the same when run as non\sphinxhyphen{}root user:

\begin{sphinxVerbatim}[commandchars=\\\{\}]
   Device Mount                          (c a) type
/dev/sda1 /efi                           (\(\pmb{\checkmark}\) \(\pmb{\checkmark}\)) EFI
/dev/sda5 /boot                          (  \(\pmb{\checkmark}\)) XBOOTLDR
\end{sphinxVerbatim}
\end{quote}
\begin{itemize}
\item {} 
\sphinxAtStartPar
Example 3

\sphinxAtStartPar
Two ESP and two XBOOTLDR:

\begin{sphinxVerbatim}[commandchars=\\\{\}]
        Device Mount                          (c a) type
/dev/nvme0n1p1 /efi                           (\(\pmb{\checkmark}\) \(\pmb{\checkmark}\)) EFI
     /dev/sdg1 /mnt/root1/efi                 (   ) EFI
/dev/nvme0n1p2 /boot                          (\(\pmb{\checkmark}\) \(\pmb{\checkmark}\)) XBOOTLDR
     /dev/sdg5 /mnt/root1/boot                (   ) XBOOTLDR
\end{sphinxVerbatim}

\end{itemize}

\sphinxAtStartPar
When run as non\sphinxhyphen{}root user, any current XBOOTLDR remains unknown/unmarked.


\subsubsection{QUIRKY ESP IN NVRAM}
\label{\detokenize{_build/man-ready/sd-boot-find-boot-mounts.8:quirky-esp-in-nvram}}
\sphinxAtStartPar
It is possible that the motherboard firmware chooses a different ESP than
that defined by nvram boot order. When this happens these 2 ESP partition
do not match.

\sphinxAtStartPar
Running \sphinxstyleemphasis{bootctl} when they are mismatched will display:

\begin{sphinxVerbatim}[commandchars=\\\{\}]
\PYG{n}{WARNING}\PYG{p}{:} \PYG{n}{The} \PYG{n}{boot} \PYG{n}{loader} \PYG{n}{reports} \PYG{n}{a} \PYG{n}{different} \PYG{n}{partition} \PYG{n}{UUID} \PYG{n}{than} \PYG{n}{the} \PYG{n}{detected} \PYG{n}{ESP}\PYG{o}{.}
\end{sphinxVerbatim}

\sphinxAtStartPar
We designate partitions determinted by boot nvram as \sphinxstyleemphasis{active} and expected partitions as
as \sphinxstyleemphasis{current}.

\sphinxAtStartPar
Please note that root privilages are required to check for a current XBOOTLDR partitions.
No elevated priveleges are otherwise needed.

\sphinxAtStartPar
To manually check for any current ESP and XBOOTLDDR partitions:

\begin{sphinxVerbatim}[commandchars=\\\{\}]
\PYG{n}{bootctl} \PYG{o}{|} \PYG{n}{grep} \PYG{o}{\PYGZhy{}}\PYG{n}{E} \PYG{l+s+s1}{\PYGZsq{}}\PYG{l+s+s1}{ESP:|XBOOTLDR:}\PYG{l+s+s1}{\PYGZsq{}}
\end{sphinxVerbatim}

\sphinxAtStartPar
And to list all active ESP and XBOOTLDR partitions:

\begin{sphinxVerbatim}[commandchars=\\\{\}]
\PYG{n}{fdisk} \PYG{o}{\PYGZhy{}}\PYG{n}{l} \PYG{o}{|} \PYG{n}{grep} \PYG{o}{\PYGZhy{}}\PYG{n}{E} \PYG{l+s+s1}{\PYGZsq{}}\PYG{l+s+s1}{EFI System|Linux extended boot}\PYG{l+s+s1}{\PYGZsq{}}
\end{sphinxVerbatim}

\sphinxAtStartPar
Also useful to see boot order:

\begin{sphinxVerbatim}[commandchars=\\\{\}]
\PYG{n}{efibootmgr}
\end{sphinxVerbatim}

\sphinxAtStartPar
and the UUID and PartUUID:

\begin{sphinxVerbatim}[commandchars=\\\{\}]
\PYG{n}{lsblk} \PYG{o}{\PYGZhy{}}\PYG{n}{o} \PYG{n}{NAME}\PYG{p}{,}\PYG{n}{FSTYPE}\PYG{p}{,}\PYG{n}{MOUNTPOINTS}\PYG{p}{,}\PYG{n}{PARTUUID}\PYG{p}{,}\PYG{n}{UUID}
\end{sphinxVerbatim}
\begin{itemize}
\item {} 
\sphinxAtStartPar
Example 4

\sphinxAtStartPar
Two ESP with mismatched active and current

\begin{sphinxVerbatim}[commandchars=\\\{\}]
        Device Mount                          (c a) type
/dev/nvme0n1p1 /efi                           (\(\pmb{\checkmark}\)  ) EFI
/dev/nvme1n1p1 /man/disk2/efi                 (  \(\pmb{\checkmark}\)) EFI
\end{sphinxVerbatim}

\end{itemize}

\sphinxAtStartPar
After eveything was done to ensure everything being correct and this \sphinxstyleemphasis{quirk} did not go away
we found that removing the offending nvram variable resolved the issue. After removing
and rebooting the nvram was then updated to the correct value by the motherboard.

\sphinxAtStartPar
To remove the nvram holding the wrong value:

\begin{sphinxVerbatim}[commandchars=\\\{\}]
\PYG{n}{cd} \PYG{o}{/}\PYG{n}{sys}\PYG{o}{/}\PYG{n}{firmware}\PYG{o}{/}\PYG{n}{efi}\PYG{o}{/}\PYG{n}{efivars}\PYG{o}{/}

\PYG{n}{echo} \PYG{l+s+s2}{\PYGZdq{}}\PYG{l+s+s2}{Should show the wrong value}\PYG{l+s+s2}{\PYGZdq{}}
\PYG{n}{cat} \PYG{n}{LoaderDevicePartUUID}\PYG{o}{\PYGZhy{}}\PYG{l+m+mi}{4}\PYG{n}{a67b082}\PYG{o}{\PYGZhy{}}\PYG{l+m+mi}{0}\PYG{n}{a4c}\PYG{o}{\PYGZhy{}}\PYG{l+m+mi}{41}\PYG{n}{cf}\PYG{o}{\PYGZhy{}}\PYG{n}{b6c7}\PYG{o}{\PYGZhy{}}\PYG{l+m+mi}{440}\PYG{n}{b29bb8c4f}

\PYG{n}{chattr} \PYG{o}{\PYGZhy{}}\PYG{n}{i} \PYG{n}{LoaderDevicePartUUID}\PYG{o}{\PYGZhy{}}\PYG{l+m+mi}{4}\PYG{n}{a67b082}\PYG{o}{\PYGZhy{}}\PYG{l+m+mi}{0}\PYG{n}{a4c}\PYG{o}{\PYGZhy{}}\PYG{l+m+mi}{41}\PYG{n}{cf}\PYG{o}{\PYGZhy{}}\PYG{n}{b6c7}\PYG{o}{\PYGZhy{}}\PYG{l+m+mi}{440}\PYG{n}{b29bb8c4f}
\PYG{n}{rm} \PYG{n}{LoaderDevicePartUUID}\PYG{o}{\PYGZhy{}}\PYG{l+m+mi}{4}\PYG{n}{a67b082}\PYG{o}{\PYGZhy{}}\PYG{l+m+mi}{0}\PYG{n}{a4c}\PYG{o}{\PYGZhy{}}\PYG{l+m+mi}{41}\PYG{n}{cf}\PYG{o}{\PYGZhy{}}\PYG{n}{b6c7}\PYG{o}{\PYGZhy{}}\PYG{l+m+mi}{440}\PYG{n}{b29bb8c4f}
\end{sphinxVerbatim}

\sphinxAtStartPar
Do this at your own risk!


\subsubsection{SEE ALSO}
\label{\detokenize{_build/man-ready/sd-boot-find-boot-mounts.8:see-also}}
\sphinxAtStartPar
\sphinxstylestrong{sd\sphinxhyphen{}boot\sphinxhyphen{}efi\sphinxhyphen{}tool\sphinxhyphen{}update}, \sphinxstylestrong{sd\sphinxhyphen{}boot\sphinxhyphen{}kernel\sphinxhyphen{}update}, \sphinxstylestrong{sd\sphinxhyphen{}boot\sphinxhyphen{}find\sphinxhyphen{}boot\sphinxhyphen{}mounts},
\sphinxstylestrong{kernel\sphinxhyphen{}install}

\sphinxstepscope


\section{sd\sphinxhyphen{}boot\sphinxhyphen{}kernel\sphinxhyphen{}update}
\label{\detokenize{_build/man-ready/sd-boot-kernel-update.8:sd-boot-kernel-update}}\label{\detokenize{_build/man-ready/sd-boot-kernel-update.8::doc}}

\subsection{Install (or remove) a kernel to (or from) boot partition}
\label{\detokenize{_build/man-ready/sd-boot-kernel-update.8:install-or-remove-a-kernel-to-or-from-boot-partition}}

\subsubsection{SYNOPSIS}
\label{\detokenize{_build/man-ready/sd-boot-kernel-update.8:synopsis}}
\sphinxAtStartPar
\sphinxcode{\sphinxupquote{sd\sphinxhyphen{}boot\sphinxhyphen{}kernel\sphinxhyphen{}update}}  \sphinxcode{\sphinxupquote{add}} | \sphinxcode{\sphinxupquote{remove}} | \sphinxcode{\sphinxupquote{inspect}} \sphinxcode{\sphinxupquote{<package\_name>}}


\subsubsection{DESCRIPTION}
\label{\detokenize{_build/man-ready/sd-boot-kernel-update.8:description}}
\sphinxAtStartPar
Install, remove or inspect a linux kernel in the \$BOOT partition using kernel\sphinxhyphen{}install.
The \$BOOT partition, following the notation used by kernel\sphinxhyphen{}install, is usually one of \sphinxstyleemphasis{/boot} or \sphinxstyleemphasis{/efi}.

\sphinxAtStartPar
It takes two arguments. The first argument is the operation:

\begin{sphinxVerbatim}[commandchars=\\\{\}]
\PYG{o}{*}\PYG{n}{add}\PYG{o}{*} \PYG{o}{|} \PYG{o}{*}\PYG{n}{remove}\PYG{o}{*} \PYG{o}{|} \PYG{o}{*}\PYG{n}{inspect}\PYG{o}{*}
\end{sphinxVerbatim}

\sphinxAtStartPar
followed by the package name providing the kernel.

\sphinxAtStartPar
The kernel package name must be one that is listed to be managed by sd\sphinxhyphen{}boot and
the kernel package must provide a package basename file that has the package name.

\sphinxAtStartPar
When a kernel package is installed using pacman, it  must be installed to:

\begin{sphinxVerbatim}[commandchars=\\\{\}]
\PYG{o}{/}\PYG{n}{usr}\PYG{o}{/}\PYG{n}{lib}\PYG{o}{/}\PYG{n}{modules}\PYG{o}{/}\PYG{o}{\PYGZlt{}}\PYG{n}{kernel\PYGZus{}version}\PYG{o}{\PYGZgt{}}
\end{sphinxVerbatim}

\sphinxAtStartPar
It must never install any files to /boot (or /efi). kernel\sphinxhyphen{}install (and hence sd\sphinxhyphen{}boot)
are responsible for using vmlinux from \sphinxstyleemphasis{/usr/lib/modules/<vers>}, generating the initrd
and installing them into \$BOOT.

\sphinxAtStartPar
The layout, specified in \sphinxstyleemphasis{/etc/kernel/install.conf} may be set to \sphinxstyleemphasis{bls} or \sphinxstyleemphasis{uki}.
In bls layout, the kernel and the initrd are installed separately. In addition
type \#1 boot loader entry files are generated.

\sphinxAtStartPar
In \sphinxstyleemphasis{uki} layout the kernel and initrd are combined into a single unified kernel image (uki)
and installed into the \sphinxstyleemphasis{efi}. These use type \#2 loader entries that are taken directly from the
uki install and do not use a separate loader entry file,

\sphinxAtStartPar
\sphinxstyleemphasis{uki} layout is preferable. It is simpler, does not require separate loader entry
files, and both kernel and initrd are signed since they are in one file. The \sphinxstyleemphasis{uki} file
itself is an \sphinxstyleemphasis{efi} file and can therefore be directly booted without the need of a boot manager
should that ever be needed.

\sphinxAtStartPar
sd\sphinxhyphen{}boot will only act if 2 conditions are met.
\begin{itemize}
\item {} 
\sphinxAtStartPar
the kernel package name is listed to be managed by sd\sphinxhyphen{}boot

\item {} 
\sphinxAtStartPar
there must be a special file installed by the package in the
\sphinxstyleemphasis{/usr/lib/modules/<kernel\sphinxhyphen{}version>} directory called \sphinxstyleemphasis{pkgbase\sphinxhyphen{}sdb} or \sphinxstyleemphasis{pkgbase}.
This file should contain the package name.

\end{itemize}

\sphinxAtStartPar
There is an related tool (\sphinxstyleemphasis{sd\sphinxhyphen{}boot\sphinxhyphen{}kernel\sphinxhyphen{}update\sphinxhyphen{}triggers}) used by pacman that is
called via an ALPM hook.

\sphinxAtStartPar
It is similar but takes a single argument only (add or remove) and reads the \sphinxstyleemphasis{triggers}
from stdin as provided by pacman.


\subsubsection{ARGUMENTS}
\label{\detokenize{_build/man-ready/sd-boot-kernel-update.8:arguments}}
\sphinxAtStartPar
\sphinxstylestrong{Operation}
\begin{itemize}
\item {} 
\sphinxAtStartPar
add <kernel\_package\_name>
\begin{quote}

\sphinxAtStartPar
Generate initrd and install the kernel provided in the package <kernel\_package\_name>
along with it’s initrd to \$BOOT
\end{quote}

\item {} 
\sphinxAtStartPar
remove
\begin{quote}

\sphinxAtStartPar
Remove the kernel from \$BOOT.
\end{quote}

\item {} 
\sphinxAtStartPar
inspect
\begin{quote}

\sphinxAtStartPar
Display information about an installed package.
\end{quote}

\end{itemize}

\sphinxAtStartPar
\sphinxstylestrong{Package Name}

\sphinxAtStartPar
Can be a kernel package name or the string “–all–” which means all kernels managed by sd\sphinxhyphen{}boot.

\sphinxAtStartPar
Note that \sphinxstyleemphasis{add} and \sphinxstyleemphasis{remove} updates files in \$BOOT. pacman is responsible for installing or removing the
package itself.


\subsubsection{FILES}
\label{\detokenize{_build/man-ready/sd-boot-kernel-update.8:files}}
\sphinxAtStartPar
\sphinxcode{\sphinxupquote{/etc/kerne/install.conf}}
\begin{quote}

\sphinxAtStartPar
The kernel\sphinxhyphen{}install configuration file. The default version provided by sd\sphinxhyphen{}boot uses
uki layout and dracut for the initrd generator:
\end{quote}

\begin{sphinxVerbatim}[commandchars=\\\{\}]
layout=uki
initrd\PYGZus{}generator=dracut
uki\PYGZus{}generator=ukify
\end{sphinxVerbatim}

\sphinxAtStartPar
\sphinxcode{\sphinxupquote{/etc/sd\sphinxhyphen{}boot/kernel.packages}}
\begin{quote}

\sphinxAtStartPar
The kernel package name must be listed here to permit sd\sphinxhyphen{}boot to manage installing
or removing the kernel
\end{quote}

\sphinxAtStartPar
\sphinxcode{\sphinxupquote{/usr/lib/modules/<kernel\sphinxhyphen{}version>/vmlinuz}}
\begin{quote}

\sphinxAtStartPar
This is the kernel, provided by the package, that will be installed to \$BOOT.
\end{quote}

\sphinxAtStartPar
\sphinxcode{\sphinxupquote{/usr/lib/modules/<kernel\sphinxhyphen{}version>/pkgbase\sphinxhyphen{}sdb}}
\begin{quote}

\sphinxAtStartPar
This file contains the name of the kernel package.
Arch kernels provide the file \sphinxstyleemphasis{pkgbase}.

\sphinxAtStartPar
While \sphinxstyleemphasis{sd\sphinxhyphen{}boot} uses either of these files, \sphinxstyleemphasis{pkgbase\sphinxhyphen{}sdb} is preferred
since the Arch kernel install tools act on any kernel using \sphinxstyleemphasis{pkgbase}.
Using a different filename prevents the Arch tools from installing kernels
sd\sphinxhyphen{}boot is already handling.
\end{quote}

\sphinxAtStartPar
\sphinxcode{\sphinxupquote{\$BOOT/<machine\sphinxhyphen{}id>/<package\_name>\sphinxhyphen{}<package\_version>}}
\begin{quote}

\sphinxAtStartPar
When layout is \sphinxstyleemphasis{bls} only,
Kernels and initrd fles are installed to this directory.
\end{quote}

\sphinxAtStartPar
\sphinxcode{\sphinxupquote{\$BOOT/loader/entries/<machine\sphinxhyphen{}id>\sphinxhyphen{}<package\_name>\sphinxhyphen{}<package\_version>.conf}}
\begin{quote}

\sphinxAtStartPar
When layout is \sphinxstyleemphasis{bls} only,
This is the type \#1 boot loader entry file. It has the title that is used in boot menu
along with the path to the program to be booted.
\end{quote}

\sphinxAtStartPar
\sphinxcode{\sphinxupquote{\$BOOT/EFI/Linux/<machine\sphinxhyphen{}id>\sphinxhyphen{}<kernel\_version>.efi}}
\begin{quote}

\sphinxAtStartPar
When layout is \sphinxstyleemphasis{uki} only,
the unified kernel image is saved here.
\end{quote}


\subsubsection{EXAMPLES}
\label{\detokenize{_build/man-ready/sd-boot-kernel-update.8:examples}}
\sphinxAtStartPar
To install the kernel provided by the package \sphinxstyleemphasis{linux\sphinxhyphen{}custom}:

\begin{sphinxVerbatim}[commandchars=\\\{\}]
\PYG{c+c1}{\PYGZsh{} sd\PYGZhy{}boot\PYGZhy{}kernel\PYGZhy{}update add linux\PYGZhy{}custom}
\end{sphinxVerbatim}

\sphinxAtStartPar
where linux\sphinxhyphen{}custom must be listed in the \sphinxstyleemphasis{/etc/sd\sphinxhyphen{}boot/kernel.packages}.

\sphinxAtStartPar
In addition the linux\sphinxhyphen{}custom package must provide the file:

\begin{sphinxVerbatim}[commandchars=\\\{\}]
\PYG{o}{/}\PYG{n}{usr}\PYG{o}{/}\PYG{n}{lib}\PYG{o}{/}\PYG{n}{modules}\PYG{o}{/}\PYG{o}{\PYGZlt{}}\PYG{n}{kernel}\PYG{o}{\PYGZhy{}}\PYG{n}{version}\PYG{o}{\PYGZgt{}}\PYG{o}{/}\PYG{n}{pkgbase}
\end{sphinxVerbatim}

\sphinxAtStartPar
or:

\begin{sphinxVerbatim}[commandchars=\\\{\}]
\PYG{o}{/}\PYG{n}{usr}\PYG{o}{/}\PYG{n}{lib}\PYG{o}{/}\PYG{n}{modules}\PYG{o}{/}\PYG{o}{\PYGZlt{}}\PYG{n}{kernel}\PYG{o}{\PYGZhy{}}\PYG{n}{version}\PYG{o}{\PYGZgt{}}\PYG{o}{/}\PYG{n}{pkgbase}\PYG{o}{\PYGZhy{}}\PYG{n}{sdb}
\end{sphinxVerbatim}

\sphinxAtStartPar
that contains one line:

\begin{sphinxVerbatim}[commandchars=\\\{\}]
\PYG{n}{linux}\PYG{o}{\PYGZhy{}}\PYG{n}{custom}
\end{sphinxVerbatim}

\sphinxAtStartPar
To inspect all kernels managed by sd\sphinxhyphen{}boot:

\begin{sphinxVerbatim}[commandchars=\\\{\}]
\PYG{n}{sd}\PYG{o}{\PYGZhy{}}\PYG{n}{boot}\PYG{o}{\PYGZhy{}}\PYG{n}{kernel}\PYG{o}{\PYGZhy{}}\PYG{n}{update} \PYG{n}{inspect} \PYG{l+s+s2}{\PYGZdq{}}\PYG{l+s+s2}{\PYGZhy{}\PYGZhy{}all\PYGZhy{}\PYGZhy{}}\PYG{l+s+s2}{\PYGZdq{}}
\end{sphinxVerbatim}

\sphinxAtStartPar
To update all kernels managed by sd\sphinxhyphen{}boot:

\begin{sphinxVerbatim}[commandchars=\\\{\}]
\PYG{n}{sd}\PYG{o}{\PYGZhy{}}\PYG{n}{boot}\PYG{o}{\PYGZhy{}}\PYG{n}{kernel}\PYG{o}{\PYGZhy{}}\PYG{n}{update} \PYG{n}{add} \PYG{l+s+s2}{\PYGZdq{}}\PYG{l+s+s2}{\PYGZhy{}\PYGZhy{}all\PYGZhy{}\PYGZhy{}}\PYG{l+s+s2}{\PYGZdq{}}
\end{sphinxVerbatim}



\renewcommand{\indexname}{Index}
\printindex
\end{document}